\documentclass[11pt]{article}
\usepackage[margin=1.1in]{geometry}
\usepackage{amsmath,amssymb,amsthm,mathtools}
\usepackage{booktabs,graphicx,enumitem,url,xcolor}
\usepackage[colorlinks=true,linkcolor=blue!50!black,citecolor=blue!50!black,urlcolor=blue!50!black]{hyperref}
\newcommand{\E}{\mathbb{E}}
\newcommand{\PP}{\mathbb{P}}
\newcommand{\R}{\mathbb{R}}
\DeclareMathOperator{\Var}{Var}
\DeclareMathOperator{\Cov}{Cov}
\DeclareMathOperator{\Leb}{Leb}
\newcommand{\eqd}{\overset{d}{=}}
\newcommand{\one}{\mathbf{1}}
\newcommand{\Imm}{\mathcal{I}}
\newcommand{\Fc}{\mathcal{F}}
\newcommand{\Lc}{\mathcal{L}}
\newcommand{\hm}{\mathfrak{h}}
\newcommand{\Ac}{\mathcal{A}}
\newcommand{\Sc}{\mathsf{S}}
\theoremstyle{plain}
\newtheorem{theorem}{Theorem}[section]
\newtheorem{proposition}[theorem]{Proposition}
\newtheorem{lemma}[theorem]{Lemma}
\newtheorem{corollary}[theorem]{Corollary}
\newtheorem{claim}[theorem]{Claim}
\theoremstyle{definition}
\newtheorem{definition}[theorem]{Definition}
\newtheorem{ansatz}[theorem]{Ansatz}
\newtheorem{convention}[theorem]{Convention}
\newtheorem{constraint}{Constraint}
\newtheorem{principle}{Principle}
\newtheorem{example}[theorem]{Example}
\newtheorem{premise}{Premise}
\newtheorem{objection}[theorem]{Objection}
\renewcommand{\thepremise}{$\Phi$\arabic{premise}}
\theoremstyle{remark}
\newtheorem{remark}[theorem]{Remark}
\title{The Route World:\\ \large A Golden-Reasoning Hypothesis of Biospheric Natural History}
\author{Ajax Benander\thanks{Email:
\href{mailto:akb9408@rit.edu}{akb9408@rit.edu}. This paper states and
develops a hypothesis; the theorem it rests on, that a record drawn under
a law conditioned on a future target has realized attributes that accord
with the target, is proved in a companion \cite{TargetedB1}, and the
survivorship model in another \cite{SurvA}. Both are cited and restated,
neither presupposed, and the hypothesis is carried throughout at a
hypothesis's grade. Drafting collaboration: Anthropic Claude (Opus~4.x,
Fable~5), under the author's direction.}}
\date{September 2026 --- Working Draft, rev.~5}
\begin{document}
\maketitle

\begin{abstract}
This paper states a hypothesis about the record of natural history and
develops what would follow from it. The hypothesis, the Route World, is
that the record is drawn under the survivorship-conditioned law of a
specific target, the Golden Point, the state past which a biosphere's
external hazard is driven to zero; the theorem that such a draw produces
realized attributes according with its target is proved in a companion
and imported, and this paper's claim is the between-model one that the
theorem does not license, that our record is drawn so. The claim in one
sentence: under the stated counterfactual-weight and target-requirement
assumptions, the record of natural history yields a Bayes factor
favoring conditioning on the biosphere's indefinite survival over
conditioning on the existence of observers, prospectively estimating
that factor is the empirical test, and every consequence drawn from the
hypothesis is carried at the grade of a hypothesis with stated tests. Four things are developed.
The target is constructed rather than stipulated: the Departure Point as
the least fixed point of an implication cascade, and the Continuity
Principle that makes one target bear on a Mesozoic mammal and on a launch
vehicle, without which no conditioning could run across the record. The
two arguments at the record's two grains, human history and natural
history, are stated as model comparisons, with the human-history leg's
flukes read as within-model diagnostics and the natural-history leg's
counterfactual weight bounded and calibrated against once-only
attributes. The retrocausal instrument is stated as an instance of the
imported theorem, with the finale paradox it dissolves and the investment
reading it licenses. And the Route World's discriminating quantities are
named as the information flux and a provisioning base rate, with the
flagship test stated prospectively: a predeclared attribute family, a
baseline, a target relation, a stopping rule, and a multiplicity budget,
declared today and scored on what follows. The retrospective cases are
labeled exploratory. What is refused is stated once: no agent, no
certification, no occupant.
\end{abstract}

\section{Introduction}\label{sec:intro}

A companion proves a theorem about records and targets: if the realized
member of a binary partition of the past is drawn under a law conditioned
on a future event, then the realized member typically accords with that
event, to within an exceedance the draw prices, and the statement is
indifferent to whether the event lies in the past or the future
\cite{TargetedB1}. The theorem is a conditional. It says what a record
looks like under a future-targeted draw. It does not say that any record
is drawn so, and it is careful to say that it does not. This paper takes
the step the theorem declines, and takes it as a hypothesis. It proposes
that the record of natural history is drawn under conditioning on a
specific future target, the indefinite survival of the biosphere, and it
asks what follows, how the proposal could be wrong, and what would show
it.

\emph{The claim, in one sentence.} Under the stated
counterfactual-weight and target-requirement assumptions, the record of
natural history yields a Bayes factor favoring conditioning on the
biosphere's indefinite survival over conditioning on the existence of
observers; prospectively estimating that factor, with both probabilities
fixed independently of the record they are scored on, is the empirical
test; and every consequence drawn from the hypothesis carries a stated
test.

\emph{What is imported.} The theorem companion supplies future-targeted
accordance, its prior-bound corollary, the conjunction form under
submultiplicativity, the multiplicity guard, the horizon-indexed
specialization to survival targets, the identification of the Bayes factor
with the Doob density, the information flux, and the limit on
certification \cite{TargetedB1}. The survivorship companion supplies the
model and its $Q$-process \cite{SurvA}. The Nexic companion supplies the
Accordance Principle at a general target, the Imperfection Magnitude Drop,
and the hyperblock construction that fixes what counts as a nameable
variation of ourselves \cite{Nexus}. Each is restated in
Section~\ref{sec:imports} at its grade.

\emph{What is new.} The target, constructed from the record's side
(Section~\ref{sec:target}). The hypothesis and its two arguments as model
comparisons (Section~\ref{sec:conditioning}). The natural-history leg's counterfactual weight, bounded from below by the single-flip sum (Section~\ref{sec:natural}). The Route World's discriminators and its prospective test (Sections~\ref{sec:routeworld} and \ref{sec:protocol}). And the fences (Section~\ref{sec:fences}). The retrocausal instrument's phenomenology, the finale paradox and the investment reading, is carried in the master manuscript and is not part of the hypothesis stated here.


\emph{Position in the branch.} The Golden branch is four papers in a
strict order of import, and this paper is its terminal synthesis. The
survivorship companion \cite{SurvA} is the root and owns the stochastic
dynamics and the conditioned path laws; the theorem companion
\cite{TargetedB1} owns the general result that a record drawn under a
target-conditioned law has realized attributes according with the
target, and the model-comparison and multiplicity machinery, and its
core theorem depends on nothing Golden; the technoanatomy companion
\cite{TechnoC} owns the mechanism, the prerequisite schedule through
which a lineage under survivorship conditioning is read as pushed, and
works with the survivorship companion's generic viability target
without assuming anything stated here. This paper imports all three and
owns the hypothesis that one particular target, long-run biospheric
survival, explains our record, together with its tests; it cites the
technoanatomy companion as a possible mechanism, and no upstream paper
cites this one, so the citation graph runs the same way as the logic.

\emph{What is refused.} No agent behind the conditioning. No
certification of the favorable side, which the imported theorem forbids
at any finite horizon. No occupant of the target's arrival. And no
inference from the theorem to the hypothesis: the theorem is within-model,
the hypothesis is between-model, and the two are never run together.

\emph{Grades.} [P] proved here or in a companion; [S] scaling-level or
structural, argued without proof; [$\Phi$] a stipulative commitment stated
as such; [O] open, with its specification. The hypothesis itself is [S]
throughout and is never upgraded by use.

\section{Vocabulary}\label{sec:vocab}
\begin{itemize}[itemsep=2pt]
\item \textbf{Departure Point.} The irreversible loss of a biosphere's
future, constructed as the least fixed point of an implication cascade.
\item \textbf{Golden Point.} Its complement; the state past which the
external hazard is driven to zero; the target.
\item \textbf{Aethus.} An observer-relative information state: stated
attributes, the unstated blank \cite{Aethus}.
\item \textbf{Cosmic Nexus / Golden Nexus.} The occurrence-frequency
weighting of attributes, and the same weighting under the target.
\item \textbf{Accordance.} The imported theorem: under a target-conditioned
draw, a realized attribute accords with the target.
\item \textbf{Past-Alignment.} The hypothesis that the record is drawn under
the Golden target; local alignment is the tested form, strong alignment
the global one.
\item \textbf{Fluke.} A route-deleting realized attribute of small prior.
\item \textbf{Stockpile.} A predeclared, submultiplicative family of flukes,
read as one attribute.
\item \textbf{Task-object.} The individuation by functional role that lets a
capacity register as missing in a history where it does not occur.
\item \textbf{The ceiling.} No event of bounded duration perfectly separates
the conditioned law from the unconditioned one.
\item \textbf{Information flux.} The rate of the relative entropy between
the two laws; the Route World's discriminating statistic.
\item \textbf{Route World.} The hypothesis: the record as the past-cone
restriction of the survivorship-conditioned law.
\end{itemize}

\section{Imported results}\label{sec:imports}

\begin{definition}[Aethus, the draw, and the general Accordance Principle; from \cite{Nexus}]
\label{imp:accordance}
An Aethus is a set of stated attributes with the unstated blank. The draw
is the postulate that at each open attribute the realized member is
distributed by the parent-conditional weights of the law in force. The
Accordance Principle, stated at a general conditioning target in a
tense-free containment chain, bounds the ratio of the target's conditional
probabilities across a realized attribute and its complement by the
attribute's odds, with $\kappa$-fold exceedances confined to probability
$1/(1+\kappa)$ for $\kappa>1$.
\end{definition}

\begin{theorem}[Future-targeted accordance; \textbf{[P]} in \cite{TargetedB1}]
\label{imp:targeted}
Let $\{A,A'\}$ be a binary partition of the past and $B$ a target with
$\PP(B)>0$. If the realized member $A^\ast$ is drawn under
$\PP_B=\PP(\cdot\mid B)$, then for every $\kappa>1$,
$\PP_B[\PP(B\mid A^{\ast\prime})/\PP(B\mid A^\ast)\ge\kappa\,\PP(A^\ast)/\PP(A^{\ast\prime})]\le1/(1+\kappa)$.
For a survival target in the critical regime the target is the horizon
family $\{T>t\}$ and the $Q$-process its limit. The theorem is
within-model; the inference from accordance to the draw is between-model
and is the Bayes factor $Q(E)/\PP(E)$, which at finite information is the
Doob density.
\end{theorem}

\begin{definition}[Imperfection Magnitude Drop, and the hyperblock; from \cite{Nexus}]
\label{imp:imd}
An Imperfection Magnitude Drop is an alternative natural history whose
joint probability with the target is drastically minute but nonzero,
reached as a targeted perturbation of one's own Aethus. The biological
hyperblock cosmos is the canvas on which any planet is representable as
an alternate history of the Earth by varying a stated set of its
attributes; it fixes the class of nameable variations of ourselves.
\end{definition}

\begin{definition}[The model, and the Golden Point in it; from \cite{SurvA}]
\label{imp:goldenpoint}
The total hazard is $h_t=B_t\varsigma_b(x_t)+\Gamma e^{-\rho x_t}$, an
external background muffled by accumulated capacity plus an internal term
the same capacity feeds or starves by the sign of $\rho$; Departure is the
killing time, viability the event that the total hazard ever accrued stays
finite, and the Golden Point the state past which the external channel is
driven to zero and viability obtains.
\end{definition}

\begin{theorem}[Two conditions for viability; \textbf{[P]} in \cite{SurvA}]
\label{imp:twocond}
Viability is the indicator of $\rho>0$ times a piecewise driver-viability
functional; the two factors are logically distinct necessary conditions,
the first a sign and the second silent on it.
\end{theorem}

\begin{theorem}[The ceiling; \textbf{[P]} in \cite{SurvA,TargetedB1}]
\label{imp:ceiling}
In the critical regime the survivorship-conditioned law is locally
absolutely continuous with respect to the unconditioned law on survivors
and globally singular; no event of bounded duration perfectly separates
the two, so certification of which law a record is under is unavailable
to any finite sample, while statistical discrimination with error remains
possible and is measured by the Bayes factor.
\end{theorem}

\begin{definition}[Task-object; from the Aethic companion \cite{Aethus}]
\label{imp:taskobject}
A task-object is individuated by the task it performs in a history rather
than by its material identity, so that a capacity is comparable across
histories in which it does and does not occur; perturbation arguments
compare a Task-Earth and a task-lineage, not an Earth and a lineage.
\end{definition}

\section{The target: the Departure Point, the Golden Point, and continuity}\label{sec:target}

\subsection{Choosing the target: continuity and non-privilege}\label{gc:sub:continuity}

Before the argument for Golden-conditioning, a prior question: why the
Golden Point rather than some more familiar target? ``Technologically
advanced civilization'' is the obvious candidate and is the wrong one,
for two reasons that compound. It is not well-defined. Any threshold one names, writing, metallurgy,
industry, spaceflight, is a stipulation, and the arbitrariness is not
cosmetic: every rarity verdict issued relative to such a target inherits
it.

And it is \emph{discontinuous with natural history}. A conditioning
target must be applicable across the whole record if the record is to be
read as one object. ``Civilization'' applies to the last few millennia
and to nothing before, so a framework built on it cannot run a single
conditioning across both the end-Cretaceous impact and the Industrial
Revolution. It must instead splice two analyses at a joint it has no
principled way to locate.

\begin{definition}[Departure Point]\label{gc:def:departure}
The \emph{Departure Point} is the event past which \emph{every}
continuation fails to advance --- the first moment at which all
scenarios lead to a failure to advance past that point. The modal
quantifier is the definition and not an embellishment of it: what is
being named is the closure of possibility, not the observation of
non-progress. \emph{This is what makes it the dual of the Golden Point.} That was
defined as the point past which a biosphere is too capable at self
preservation to regress past it --- which, stated with the quantifiers
exposed, is the point past which \emph{some} available course has
\emph{every} continuation avoiding foreclosure. The asymmetry between
the two is deliberate: departure is $\forall\exists$ --- every course
admits a foreclosing continuation --- and Goldenness is
$\exists\forall$ --- some course admits none.

Two $\forall$-statements could not be complements of one another;
$\exists\forall$ and $\forall\exists$ can be, and
\S\ref{gc:subsub:cascade} shows that these two are, by exhibiting them
as the complementary fixed points of De Morgan-dual operators. The
quantifier asymmetry is therefore not a wrinkle in the duality but its
mechanism. The pair is exhaustive on definite states, and an Aethus carries weight
across both. \S\ref{gc:subsub:grading} makes that precise, since it is
best stated once the grading is in hand. The reading this forces bears stating, since a looser gloss
would lose it. Goldenness is not the guarantee that nothing can go
wrong; it is the possession of a \emph{course} along which nothing
does, whatever the environment answers. It is a claim about available
strategy, not about the absence of hazard --- which is why a Golden
biosphere still faces a universe that would destroy it and is Golden
anyway.

Both are \emph{points} for the same reason: irreversibility. A state
one may still leave is not a point but a phase, and neither concept is
about phases. Two consequences of stating it this way. \emph{Biological death is a mechanism, not the definition.} Death is one
way a trajectory is foreclosed, and it happens to be the most visible
one; irreversible loss of the capacity to advance is another, and the
definition is indifferent between them. What is extracted is the
\emph{phenomenology of foreclosure}, taken apart from the particular
mechanism of biosphere-death in which it is most legible --- and once
extracted it applies wherever there is a trajectory to foreclose,
which is why the notion carries across the scales
Principle~\ref{gc:prin:continuity} requires.

\emph{Non-advancement alone does not suffice.} A biosphere that fails
to progress for an age has not thereby reached its Departure Point; it
has had a lull, and a lull is reversible. Only when no continuation
remains that advances has the point been passed. This is the
distinction that the phrase ``failure to advance'' elides and that the
quantifier restores. Two consequences for how the model of Section~\ref{sec:imports} should be read.
The hazard $h^{\mathrm{tot}}_t$ of the survivorship companion \cite{SurvA} is the hazard
\emph{of foreclosure}, not of death; the muffling exponent $x$ is
accumulated against that; and every Imperfection Magnitude Drop is a
movement toward it. And the analytic form of the modal claim is \emph{pathwise}
divergence, $\Lambda_\infty=\infty$ along every continuation, which
the trapped regime satisfies outright
(the survivorship companion \cite{SurvA}(i)), of which
$\PP(\Lambda_\infty<\infty)=0$ is the measure shadow; the two come
apart exactly where the broader Golden program, set aside here separates
them, since the critical regime carries the shadow while advancing
continuations survive at measure zero. This is why the
trapped regime of the survivorship companion \cite{SurvA}(iii) has
departed: not because a biosphere halting at $x_\ast$ has stopped, but
because with $\rho<0$ the viability probability is identically zero, so
no continuation available to it advances. Halting is not the departure;
having nowhere left to go is.

The term is used rather than ``extinction'' deliberately, and the
reason is the same one that governs the choice of the Golden Point
itself. ``Extinction'' is a biological category, applying to taxa and
carrying the vocabulary of one discipline; the Departure Point must
apply at every scale the record contains, to a Mesozoic lineage as
readily as to a biosphere tomorrow, and must name an event in the
Aethic block universe rather than a fact about organisms. It is
therefore defined at the same primitive level as the Golden Point,
which is what allows Principle~\ref{gc:prin:continuity} to be stated at
all.

\emph{Internal and external.} The Departure Point divides by the
locus of its cause. An \emph{internal} Departure Point is brought about
by factors within the biosphere's own homeworld; an \emph{external} one
is inflicted from outside, and is roughly constant in magnitude over
long timescales because the cosmos is roughly constant in makeup. That
division is the ontological origin of the two channels of
the survivorship companion \cite{SurvA}: the crest channel carries the external hazard and
is bounded, the trough channel carries the internal and is not.
Technoanatomy (the survivorship companion \cite{SurvA}) is then the statement
of how the second responds to capacity acquired against the
first.\footnote{\emph{Provenance.} The Departure Point, its
internal/external division, the constancy of the external threat, and
the observation that a biosphere may come to draw energy out of itself
are all set out in our \emph{Golden Science} of August 2022
\cite{GoldenScience22}, some years before the present formalism and
before the drafting collaboration. That document also states the Golden Point as the point
past which a biosphere cannot regress, derives a positive immortality
probability from an exponentially decaying hazard, argues for
``biosphere'' over ``civilization'' as the unit, and names the Trace
Connection as an open problem --- which \S\ref{gc:sub:pastalign} is the
attempt to close. The formalism here should be read as supplying
machinery for commitments already made there, not as generating
them.}
\end{definition}

\subsubsection{The implication cascade: departure as a least fixed point}
\label{gc:subsub:cascade}

The construction imports the companion's fixed-point result rather
than re-deriving it, and reads the descendant relation $\sqsubseteq$
as the companion now reads it, transitively; nothing below depends on
generation counts or on finite branching. The passage from the model's
verdict to the modal set is a bridge, not an identity: $\Lambda_\infty=\infty$ gives almost-sure eventual Cox
killing, while membership in the doom set $\mu G$ means every
admissible continuation is doomed, and the two coincide only under the
identification of zero-weight continuations with absent ones, which
the framework adopts and states here rather than passing from
$\PP(T=\infty)=0$ to $A\in\mu G$ without it.

Definition~\ref{gc:def:departure} carries a recursion that is worth
making explicit, because it is already available in the companion
framework in finished form and because it changes what the model's
hazard is a hazard \emph{of}. The internal and external hazards of the survivorship companion \cite{SurvA} \emph{imply}
departure; they are not exhaustive of it. A biosphere may reach a
configuration from which no route avoids those hazards, and by the
modal definition that configuration \emph{is} a Departure Point,
attained before either hazard has fired. The statement ``no matter what
you do, you reach a Departure Point'' is itself a Departure Point ---
and since that statement can again be true one step earlier, the
implication propagates.

\begin{remark}[This is the third postulate's operator, imported]
\label{gc:rem:thirdpostulate}
The companion's third postulate supplies exactly this recursion
\cite{Aethus}, and its result is imported as stated there. Over
Aethae ordered by the descendant relation $\sqsubseteq$, with
$\mathrm{Base}$ collecting the non-recursive modes, the safety operator
and its dual are
\begin{align*}
F(Y)&=\bigl\{A\notin\mathrm{Base}:\ \exists B\sqsubseteq A,\
\forall C\sqsubseteq B,\ C\in Y\bigr\},\\
G(X)&=\mathrm{Base}\cup\bigl\{A:\ \forall B\sqsubseteq A,\
\exists C\sqsubseteq B,\ C\in X\bigr\},
\end{align*}
both monotone, so that by Knaster--Tarski the greatest fixed point
$\nu F$ and the least fixed point $\mu G$ exist for arbitrary branching
and are complements of one another; the companion's seeded convention
fixes $\mathrm{Base}$, and under it $\mathcal D=\mu G=G(\mathrm{Base})$
closes. Generation counts, finite branching, and closure ordinals play
no role in the imported statement and none here.

Read $\mathrm{Base}$ as departure by internal or external hazard
--- departure-sufficient occurrences in the sense of
Definition~\ref{imp:goldenpoint}, with nothing about death presupposed.
Then $G$ is the doom operator in the sense above, every route has a
doomed continuation, and the Departure Point is
\begin{equation}\label{gc:eq:departurefix}
\mathcal{D}\;=\;\mu G\;\supsetneq\;\mathrm{Base}.
\end{equation}
Three consequences, each already carried by the companion's
machinery rather than assumed here.

\emph{The implication is one-directional by design.} The companion
states the third postulate as an implication and not a biconditional,
which is precisely the containment in \eqref{gc:eq:departurefix}: the
hazards suffice for departure and are not necessary to it. \emph{The Golden Point is an event horizon in a precise sense.} It is
the state past which the capacity to trigger departure is closed off,
which is $\nu F$'s $\exists\forall$ said in the other direction, and
the Departure Point is its dual, the state at which every remaining
possibility retains that capacity; the parallel is recorded rather than
used.

\emph{The Golden/Departure duality is a theorem, not a construction.}
The Golden Point is $\nu F$, there exists a route all of whose
continuations remain safe, and the Departure Point is $\mu G$. Since
$G=\neg F\neg$, the two fixed points are complements outright, so the
duality asserted in Definition~\ref{gc:def:departure} need not be
stipulated: it is inherited, and the pair is exhaustive. \emph{The complementarity is a fact about definite states.} The
fixed-point analysis lives on the Boolean lattice of definite states,
where the result above is exact. An Aethus is a weighted superposition
of such states, and what that implies for the categories is taken up in
\S\ref{gc:subsub:grading}.

\emph{Human extinction sits in the difference
$\mu G\setminus\mathrm{Base}$.} It is not itself a biosphere-level
departure by internal or external hazard,
since the biosphere persists; but with no muffling agent remaining,
every continuation reaches $\mathrm{Base}$
(the broader Golden program, set aside here). It therefore enters at the
first step of the cascade. The classical reading, that human
extinction is a Departure Point because it is only a matter of time, is the $n=1$ case of the recursion, and \eqref{gc:eq:departurefix} is
what licenses it.
\end{remark}

\begin{remark}[The model's hazard is a lower bound]
\label{gc:rem:hazardlower}
The consequence for Section~\ref{sec:imports} should be stated plainly, since it
qualifies every rate in this work. The hazard $h^{\mathrm{tot}}_t$ of
the survivorship companion \cite{SurvA} is the hazard of $\mathrm{Base}$: it prices the
internal and external mechanisms and nothing else. The hazard of
\emph{departure} is the hazard of $\mu G$, which by
\eqref{gc:eq:departurefix} is at least as large and in general strictly
larger --- because a biosphere departs when its fate is sealed, not
when the sealing is executed, and the interval between those can be
long.

So the model's timescales are conservative in a specific direction:
where it reports a departure time $T$, the Departure Point proper may
have been passed considerably earlier, and the survival curves of
Section~\ref{sec:imports} are upper bounds on the probability of not
having departed. Nothing in the asymptotic analysis is disturbed by
this, a cascade that fires earlier by a bounded lead time leaves
exponents untouched, but the quantitative readings inherit it, and
the anatomy urgency of the survivorship companion \cite{SurvA} is sharpened by it: the
window in which one may still act closes at the cascade's rate, not at
the hazard's. Computing that rate is open \textbf{[O]};
the schedule of \cite{TechnoC} is a first attempt at it, exhibiting the excess
as a third hazard channel whose completion weight is a least fixed point
the technoanatomy companion \cite{TechnoC}.

A second approach, complementary rather than competing, is available in
the companion and is worth recording here because it constrains what the
rate can be. A sealing-hazard is not a hazard on a trajectory but a
\emph{settlement} rate on the books: what it prices is the acquisition of
statedness by the attribute that the fate is sealed, and that is an
event of the present index whatever its referent's date --- statedness
being a property of the indexing Aethus and never of a clock reading, so
that blanks lack tense \cite{Aethus, AethusDecision}. Three consequences follow, and the first is
restrictive. The weight the current books place on a fixed
future-referring attribute is, by the framework's law of total
probability, a \emph{bounded martingale} along the progression, the
current state prices its own next resolutions, so its expected rate of
increase vanishes identically and \emph{the sealing-hazard cannot be a
drift}. What survives as a per-unit-time object is the settlement flux,
the weight mass absorbed into statedness per unit proper time, together
with the direction-blind quadratic-variation rate, which a refined
account could consult alongside it. And since a bounded martingale
converges, the current weight audits its own hazard: it decomposes into
the price of eventual settlement plus the expected never-settled residue,
so the number available at an instant is already an integral over the
flux.

That last observation has a consequence for how the open problem should
be posed, and it is the one worth carrying. The flux and the level are
not rival candidates for \emph{the} sealing-hazard; they are the two ends
of a one-parameter family, obtained by pricing settlement at a discount,
\begin{equation}\label{gc:eq:sealprice}
h_\gamma(B)\;=\;\E_B\bigl[e^{-\gamma(\rho-t)}W_\rho\bigr],
\qquad \gamma\ge0,
\end{equation}
with $\rho$ the earlier of settlement and the path's terminus. At
$\gamma\to0$ optional stopping returns the level itself; at
$\gamma\to\infty$ the exponential concentrates and $\gamma h_\gamma$
returns the instantaneous rate. And the kernel's shape is \emph{forced}
rather than chosen: requiring the pricing to commute with refinement of
the progression, the price of the children's prices returning the
parent's, makes the delay factor multiplicative across concatenated
intervals, whose bounded monotone solutions are exactly the exponentials
\cite{Aethus}.

Two things follow for this work. The exponential of the Cox construction
(the survivorship companion \cite{SurvA}) is therefore derivable from partition consistency
rather than adopted with it. And this paper already uses both ends of the
family without having named them as ends: the survivorship companion \cite{SurvA}'s
$h^{\mathrm{tot}}$ is a rate, the $\gamma\to\infty$ face; the schedule
channel's completion weight the technoanatomy companion \cite{TechnoC} is a level, the
$\gamma\to0$ face. Reading them as one family says what the open rate of
this remark actually is --- not a single number awaiting computation, but
a curve whose urgency parameter is itself a modelling commitment, and
whose mixtures over urgencies are, by Bernstein's theorem, exactly the
completely monotone kernels.

The construction is developed at individual scale in \cite{Aethus}, where
the settling attribute is lineage termination; what transfers here is the
form, by Principle~\ref{gc:prin:continuity}, and what does not transfer
is any of its quantities.
\end{remark}

\begin{proposition}[One permanent floor suffices; \textbf{[P]} given the
model]\label{gc:prop:onefloor}
Suppose some channel $h_j$ of the survivorship companion \cite{SurvA} acquires a
\emph{permanent floor}: there are $c>0$ and $t_0$ with
$h_j(t)\ge c$ for all $t\ge t_0$, no route remaining by which it falls
arbitrarily again. Then $\Lambda_\infty=\infty$, so
$\PP(\Lambda_\infty<\infty)=0$ and the state lies in $\mu G$.
\end{proposition}

\begin{proof}
The channels are non-negative, so $h^{\mathrm{tot}}\ge h_j$ pointwise,
and
\[
\Lambda_\infty=\int_0^{\infty}h^{\mathrm{tot}}_t\,dt
\;\ge\;\int_{t_0}^{\infty}h_j(t)\,dt
\;\ge\;\int_{t_0}^{\infty}c\,dt=\infty. \qedhere
\]
\end{proof}

\noindent Two things the proof does \emph{not} use are worth marking,
since both are tempting and one is false. It uses no property of the
other channels: in particular the internal channel need not
participate, and a permanent floor in the external channel alone
places the state in $\mu G$ --- which is the claim's whole point, that
one does not need the base hazard entire. And it does not pass through
``$\Lambda_\infty<\infty$ requires $h^{\mathrm{tot}}_t\to0$'', which
would be the natural route and is not available: a non-negative
function may have finite integral without tending to zero, as tall
narrow spikes of shrinking width show. The floor hypothesis makes the
detour unnecessary, since it bounds the integrand below on a set of
infinite measure directly.

\begin{example}[Human extinction as an Imperfection Magnitude Drop]
\label{gc:ex:humanIMD}
Human extinction is an Imperfection Magnitude Drop, and for reasons
that are structural rather than sentimental. The grades differ across
the parts and are marked where they fall: (A) and (B) below are claims
about historical individuation \textbf{[$\Phi$]}, their combination
with Proposition~\ref{gc:prop:onefloor} is \textbf{[P]} given the
model, and the smallness of $\varepsilon$ is imported at the
companion's grade (Remark~\ref{gc:rem:whatcarries}). Two reasons, kept
separate because they are referred to separately below:

\begin{description}\itemsep2pt
\item[\normalfont(A) \emph{The lineage is spent.}] The cascade to human
evolution ran through the human--chimpanzee last common ancestor, a
particular historical individual which no longer exists to be run
through again.
\item[\normalfont(B) \emph{The setting is spent.}] The specific sequence
of climatic and geographical events that brought that lineage down to
the grasslands has been cashed out, the world now standing at a
different stage.
\end{description}

\noindent The pathway is therefore not merely improbable on repetition
--- it is \emph{unrepeatable}, its preconditions spent. Two claims must be kept apart here, since the argument moves between
them and they carry different values. That \emph{the human route
recurs} is not improbable but unavailable: (A) and (B) name stated
attributes that are gone, so this contributes nothing at all to what
follows. That \emph{the Golden Point is reached without} (A) and (B) is
a different proposition, and it is minute rather than zero --- an
Imperfection Magnitude Drop, in the exact sense the next paragraphs
derive. Everything positive that remains after human extinction belongs
to the second; the first is simply closed.

\emph{The derivation, run on (A) and (B).} An Imperfection Magnitude
Drop is a targeted perturbation away from one's own Aethus, yielding a
natural history whose success-intersection is drastically minute but
never quite zero \cite[\S IMD]{Nexus}. Human extinction complements (A)
and (B): the last common ancestor is no longer a stated attribute
available to be run through, and the climatic-geographical setting has
been cashed out. Any Aethus continuing from that state therefore
carries those two attributes complemented.

Two established results then apply, in an order that matters. The
High-Contrast Earth Hypothesis fixes the \emph{magnitude}: complementing
an attribute correlated with the remainder of the Aethus displaces
$P(H\mid\cdot)$ by many orders of magnitude, to an effectively random
order --- the hypothesis asserting the size of the gap and not its sign
\cite{Nexus}. The Accordance Principle then fixes the \emph{direction},
placing ours on the greater side. The perturbed Aethae are accordingly
drastically minute, which is what it is to be an Imperfection Magnitude
Drop.

\emph{Why the deep-time reply does not reach it.} ``Re-evolving humans
is unlikely'' invites the answer that deep time makes unlikely things
happen, and a post-extinction biosphere has deep time in abundance. That
answer would land against a claim about rates. It does not land here,
because the displacement is a property of the \emph{perturbation} rather
than of a waiting period: complementing a correlated attribute
re-randomizes the success-intersection wherever in time the continuation
is read, and elapsed time supplies no mechanism returning it to the
unperturbed value. The claim is thus made within probability and not
outside it, as the Disaster-Difficulty Unification Principle requires
--- the two sit on the severity spectrum wherever their probability puts
them, and what distinguishes them is not their kind but that their
probability does not accumulate.

\subsubsection{Alternative lineages, by disjointness}
\label{gc:subsub:disjointness}

The objection that remains is the one about successors: humans will not
return, but some other lineage, another primate, or the corvid or
cetacean lineages in which complex cognition has arisen independently, might find its own way to the Golden Point. It is answered by the
derivation already given, and by a single observation about what such a
lineage is.

\begin{lemma}[Disjointness]\label{gc:prop:disjoint}
Let $L$ be any extant lineage other than our own, and let $A_L$ be the
Aethus of a natural history in which $L$ reaches the Golden Point. The
attributes of $L$ mediating its route are \emph{disjoint} from those
mediating (A) and (B): $L$ did not descend through the human--chimpanzee
last common ancestor, and was not the lineage driven into the African
grassland setting. The disjointness is definitional, what makes $L$ an
alternative lineage rather than ours, and it establishes exactly this:
$A_L$ carries (A) and (B) complemented, so $L$ is not the human route.
\end{lemma}

\begin{corollary}[Conditional Imperfection Magnitude Drop;
\textbf{[S]} under the landscape hypothesis]\label{gc:cor:condimd}
Under the additional hypothesis that sufficiently route-distant
alternatives are Golden-suppressed, the cumulative-suppression package
of the High-Contrast Earth Hypothesis and the Optimal-Earth companion
\cite{Ben26Optimal}, each such $A_L$ is an Imperfection Magnitude Drop
relative to our own Aethus. Without that hypothesis nothing of the kind
follows from the lemma, since a different lineage could reach the same
functional capacity by a different route, and the task-object
individuation of the task-object construction (Definition~\ref{imp:taskobject}) abstracts from material identity to
functional role by design.
\end{corollary}

Two clarifications, since the corollary is easy to over-read.

\emph{It does not say a successor is impossible.} An Imperfection
Magnitude Drop is minute and never zero, so $A_L$ retains positive
weight for every $L$. What is claimed is that the weight is displaced
by the contrast, not that it is annihilated --- and the residual is
precisely the $\varepsilon$ of \S\ref{gc:subsub:grading}. \emph{It does not depend on which lineages are named.} The primates,
corvids and cetaceans are referents rather than premises: they say what
the disjoint attributes belong to. The proposition would hold for a
lineage nobody has thought to nominate, since disjointness from (A) and
(B) is a consequence of being a different lineage and not a fact
discovered about any particular one.

\subsubsection{The conclusion}\label{gc:subsub:imdconclusion}

The assembly runs under the landscape hypothesis of
Corollary~\ref{gc:cor:condimd}, and carries its grade.

The pieces now assemble, and the last step is worth running in full
rather than by reference, since it is where a claim about \emph{lineages}
becomes a claim about the \emph{Departure Point} and the two are not
obviously the same kind of thing. (A) and (B) remove the human route;
Proposition~\ref{gc:prop:disjoint} makes every alternative route an
Imperfection Magnitude Drop. Consider then any branch carrying no new
muffling agent. On such a branch no route remains by which one
re-arises, that is what the two previous results jointly establish, and a channel whose comparisonn depends on an agent that no longer
exists and cannot re-arise is a channel that will never again fall.
Fix any $c$ below its prevailing value and any $t_0$ after the loss:
$h^{\mathrm{ext}}(t)\ge c$ for all $t\ge t_0$, which is precisely a
\emph{permanent floor} in the sense of
Proposition~\ref{gc:prop:onefloor}.

That proposition then closes the argument in two lines. Channels are
non-negative, so $h^{\mathrm{tot}}\ge h^{\mathrm{ext}}$ pointwise, and
\[
\Lambda_\infty\;\ge\;\int_{t_0}^{\infty}h^{\mathrm{ext}}(t)\,dt
\;\ge\;\int_{t_0}^{\infty}c\,dt\;=\;\infty ,
\]
so $\PP(\Lambda_\infty<\infty)=0$. By
Definition~\ref{gc:def:departure} that is the analytic form of the
modal claim, no continuation advances, and the state therefore
lies in $\mu G$. One such channel suffices, and no property of the
other channel, of $\rho$, or of the driver is used anywhere in the
step: a single floor anywhere in the total hazard is enough, whatever
the rest of the model is doing.

It is worth marking what this does \emph{not} require. It does not
require that the floor be high, only that it be positive and permanent;
a biosphere left at a very low but irreducible external hazard departs
just as surely as one left at a high one, merely later. And it does not
require extinction to be the mechanism --- what carries the step is the
loss of the \emph{agent}, and any event removing the capacity to muffle
without removing the biosphere would carry it identically.
On every such branch, therefore, human extinction is a Departure Point
\emph{outright} --- in the modal sense of
Definition~\ref{gc:def:departure}, and not merely a step toward one.

Across the full Aethic weight the statement is graded rather than
outright. Positive weight survives on the branches
Proposition~\ref{gc:prop:disjoint} prices, those on which some
alternative lineage reaches the Golden Point, and $\varepsilon$ is
that weight. It is positive because an Imperfection Magnitude Drop is
minute and never zero, and small because each such branch takes the
displacement of the High-Contrast Earth Hypothesis
(Remark~\ref{gc:rem:whatcarries}). Its formal statement is
\S\ref{gc:subsub:grading}.

\begin{remark}[What carries the argument, and what would break it]
\label{gc:rem:whatcarries}
The conclusion closes an objection that would otherwise stand against
the broader Golden program, set aside here, that a successor lineage might re-evolve the muffling
capacity, so the accumulation is inherited rather than forfeited, and
it closes it with the same machinery that establishes the drop, not
with an additional argument. What is contributed here is only the
identification of (A) and (B) as the attributes perturbed, together
with the observation that any alternative lineage is disjoint from them
(Proposition~\ref{gc:prop:disjoint}); the displacement is then supplied
by the High-Contrast Earth Hypothesis and its direction by the
Accordance Principle, both imported at the companion's grade
\cite{Nexus, Ben26Optimal}.

Three things would break it, and they attack different components.
That (A) or (B) is misidentified, that the last common ancestor or
the climatic sequence is not in fact a stated attribute of the Aethus in
the relevant sense, is an attack on this section's own contribution.
That the contrast is not vast, or that the interrogated attributes fall
outside the correlated class the Hypothesis is scoped to, is an attack
on the imported magnitude. And that the direction runs the other way is
an attack on the imported Accordance step. The first is where a critic
should aim if the target is the present argument rather than the
framework it runs on.
\end{remark}
\end{example}

\subsubsection{Grading, and why two states suffice}
\label{gc:subsub:grading}

The cascade settles what a Departure Point \emph{is}. What remains is
that no realistic state satisfies it exactly, since some weight always
survives on branches that advance --- so the operative statements are
graded, and the grading attaches to weight rather than to kind.

\begin{remark}[The recovery branch: a genuine Aethic superposition,
and the grading it supports]\label{gc:rem:boltzmann}
An Imperfection Magnitude Drop does not drive the probability of later
recovery to zero, and this is definitional rather than incidental: a
drop is a perturbation whose success-intersection is drastically minute
\emph{but never quite zero} \cite{Nexus}. The weight remaining after
human extinction is therefore exactly the drop's residual weight ---
that of the perturbed Aethae in which the Golden Point is attained
regardless. Nothing further is needed to say why it is positive, and
nothing further is needed to say why it is small: the first is what an
Imperfection Magnitude Drop is, and the second is the displacement of
Example~\ref{gc:ex:humanIMD}.

That small positive weight sits in apparent tension with a Departure
Point defined by a strict modal quantifier: if \emph{some} continuation
advances, however faintly, then not every continuation fails, and the
label should not apply. \emph{The post-drop state is an Aethic superposition, and that is the
whole of the resolution.} An Imperfection Magnitude Drop does not
replace one state by another; it produces a state carrying both a
departed component and a recovery component, each genuinely stated, in
whatever proportion the physics dictates \cite{WeightAlg}. Neither
component is an artifact of description, and neither is deleted by the
other being large. The recovery branch is therefore \emph{present in
the Aethus}, not absent from it and merely improbable --- which is
precisely the condition a strict universal quantifier is sensitive to.

\emph{So the label does not attach, and grading is forced rather than
chosen.} If a component that advances is genuinely stated, then it is
false that every continuation fails, and
Definition~\ref{gc:def:departure} does not apply however small that
component's weight. What the situation calls for is not a verdict about
whether the state is departed but a measurement of how much of it is:
the superposition is graded, so the description of it must be graded
too. Nothing is rounded and nothing is discarded to reach that
conclusion --- the grading is read off the state rather than imposed on
it.

\begin{remark}[The graded object at individual scale, and a result that
transfers]\label{gc:rem:deathgrade}
The definition below has an exact individual-scale instance in the
companion, and the companion proves something about it this work does
not: for a child $B$ of a state with survival weight $s_A$, define its
\emph{death grade} $g(B):=1-s(B)/s_A$, the fraction of survival weight
the child destroys. Then the level-price is \emph{affine} in the grade,
and a graded slide of grade $g$ is equivalent at zero urgency to a
mixture placing weight $q=g$ on total settlement and $1-q$ on no change
\cite{Aethus}. Read across, an $\varepsilon$-Departure Point is a child of
death grade $1-\varepsilon$ at biosphere scale, and the affinity and the
mixture equivalence transfer unchanged --- which is worth recording
because it says the grading is linear in exactly the coordinate the
survivor's own accounting uses (the survivorship companion \cite{SurvA}), and
not in its logarithm.
\end{remark}

\begin{definition}[$\varepsilon$-Departure Point]\label{gc:def:epsdep}
A state is an \emph{$\varepsilon$-Departure Point} when the total
Aethic weight of continuations that advance is at most $\varepsilon$. A
Departure Point in the strict sense of
Definition~\ref{gc:def:departure} is the case $\varepsilon=0$.
\end{definition}

Human extinction is then an $\varepsilon$-Departure Point with
$\varepsilon$ the weight of the recovery branch --- which is to say,
in the phrase that states the situation exactly, \emph{almost precisely
a Departure Point, probabilistically}. That is a stronger claim than
the binary label would have been, not a weaker one: it carries a
number where the label carried only a verdict, and it is true where the
label was not.

This settles the question the definition deferred.

\begin{remark}[Where the variation lies]\label{gc:rem:pairsuffices}
Write $w(A)$ for the total Aethic weight that state $A$ places on
continuations which advance. The categories are $\nu F$ and $\mu G$,
complementary and exhaustive on definite states
(\S\ref{gc:subsub:cascade}). What varies across biospheres is therefore
not \emph{which} category one occupies but \emph{how weight is
distributed} across the two: a state at $0<w<1$ is a superposition of
the pair in the proportion $w:1-w$, exactly as
Remark~\ref{gc:rem:boltzmann} describes the post-drop state.

Three consequences.

\emph{The weights carry what a finer taxonomy would.} The framework
supplies Aethic weights already, and $w$ is a number where a category
would be a verdict; the finer information is in the weight and is
available without further vocabulary. \emph{The grading attaches to the state, not to the kind.}
Definition~\ref{gc:def:epsdep} grades by weight, which is the correct
place for grading to sit: an $\varepsilon$-Departure Point is a state
whose Departure-tending weight is $1-\varepsilon$.

\emph{Every actual biosphere is superposed, and this is the subject
matter.} None is secured and none has $w=0$ while advancing
continuations remain, so every case of interest carries weight on both
poles. That is also why the technoanatomy companion \cite{TechnoC}'s interpolation $\PP^{(t)}$ is
the right object: it tracks the redistribution of weight between the
two, rather than arrival at either.
\end{remark}

\begin{remark}[Where the strictness survives, and the two layers]
\label{gc:rem:twolayers}
Two things should be kept apart, because the grading does not reach
everywhere.

\emph{The model's verdicts remain strict.} Within Section~\ref{sec:imports}, the
configuration $x\equiv0$ gives $\PP(\Lambda_\infty<\infty)=0$
\emph{exactly}, and adverse anatomy gives it exactly for every
$\rho\le0$: these are identities, not small numbers. The $\varepsilon$
of Definition~\ref{gc:def:epsdep} does not live inside the model at
all. It lives at the Aethic level, where the weighting distributes over
\emph{which configuration obtains} --- and the recovery branch is
precisely the branch on which a new muffling agent arises, which the
model of a fixed biosphere does not represent. So the model says what
happens given a configuration, strictly; the Aethic weight says how
much weight each configuration carries.

\emph{This is what the two-layer algebra is for.} The companion's
construction is an unweighted layer carrying the structural facts and a
weighted layer laid over it without disturbing them
\cite{WeightAlg}. The cascade of \S\ref{gc:subsub:cascade} is a Boolean
recursion and belongs to the first: $\mu G$ is a set, computed from
implication alone. The grading of
Definition~\ref{gc:def:epsdep} belongs to the second. Neither
displaces the other. Making the cascade itself graded would require a
fixed point of an $M$-valued doom operator, which is open \textbf{[O]}
and is the natural next question this remark raises.
the technoanatomy companion \cite{TechnoC} exhibits one concrete instance of such a
fixed point --- single-type and carrying one number rather than ranging
over the lattice, but computed rather than assigned.
\end{remark}
\end{remark}

\begin{principle}[Continuity; \textbf{[$\Phi$]}]\label{gc:prin:continuity}
The struggle our ancestors underwent against the Departure Point in
natural history is continuous, philosophically and ontologically, with a
society's agency-mediated struggle against the Departure Point. They are
not analogs but the same activity under different mediation.
\end{principle}

The Golden Point is the target that respects this, and it does so
because of what kind of definition it has: \emph{survival conditioning
of its type}. That is primitive enough to apply without alteration to a
Mesozoic mammal evading theropod predation and to a modern society
avoiding its own weapons. Both lower a hazard; both prevent an
Imperfection Magnitude Drop; neither requires the vocabulary of the
other. The substance is continuous where the societal construct is not.

\begin{objection}[Agency]\label{gc:obj:agency}
A Mesozoic mammal does not know it is fighting the Departure Point.
There is no agency, no goal, no strategy --- only an animal avoiding a
predator. Calling that continuous with deliberate societal
risk-management equivocates on the word ``fighting''.
\end{objection}

\noindent The objection assumes agency is what individuates the
activity. It is not. By the task-object construction (Definition~\ref{imp:taskobject}), what individuates a task-object
is its \emph{contribution to the root}, and agency is nowhere in that
criterion --- which is precisely why task-objects were required before
this principle could be stated. What replaces agency is the mediating mechanism, and Optimation
conditions on the mechanism rather than around it. Where a society
lowers its hazard through deliberation, a lineage lowers its through
evolutionary selection; the conditioning does not care which, because
what it fixes is that the hazard be lowered. It therefore reaches
\emph{through} the mechanism to whatever contingencies the mechanism
requires in order to operate.

\begin{example}[Conditioning through the medium: the
$\alpha$-gal case; \textbf{[S]}]\label{gc:ex:medium}
Our $\alpha$-gal case exhibits the structure exactly, and its
details matter because they show the conditioning having no alternative
route \cite{Nexus}.

Some thirty million years ago the Catarrhine lineage came under
pathogen pressure of a particular kind. The pathogen presented
$\alpha$-gal, a carbohydrate our own cells also produced, so that no
immune response was available which did not attack the host. What
survived carried a mutation disabling \textit{GGTA1}, the gene
responsible for producing it --- a gene retained by every other
mammalian lineage, and lost in ours alone.\footnote{The inactivation of
\textit{GGTA1} in Old World primates is established, as is
pathogen-mediated selection as its leading explanation. The
\emph{severity} of the episode, whether it approached extinction of
the lineage, is an inference and not a finding, and the argument does
not rest on it: pressure sufficient to remove a gene universal among
mammals is all that the reasoning requires, and a milder episode
delivering that pressure would serve the example unchanged.}

Now read it under Optimation. The gene was ancient and useful, which is
why it is universal among mammals; and evolutionary selection removes a
gene if and only if there is pressure to remove it. That is the
mechanism's own condition of operation, and it is not negotiable. So
under ordinary conditions the loss was not merely improbable but close
to unavailable: selection had nothing to act on, because the gene was
not costing anything.

A conditioning requiring the loss therefore cannot act on the gene. It
must act on whatever \emph{supplies the pressure} --- and what supplied
it was the contingent existence of a pathogen presenting the very
carbohydrate in question, which converts the gene from an asset into the
one thing standing between the lineage and extinction. Optimation
reweights the pathogen's occurrence to prominence, and the loss follows
\emph{through} selection rather than around it.

This is the general shape. Optimation does not reach past evolutionary
selection to install outcomes; it reweights the contingencies that make
selection deliver them. Agency, where it exists, is one such mechanism
among others, the deliberative rather than the Darwinian medium, and the conditioning is indifferent between them.
\end{example}

\begin{remark}[What the principle buys, and its cost]
\label{gc:rem:continuitycost}
It buys the object the rest of this work analyzes. Without continuity,
Section~\ref{sec:imports}'s hazard process would model two different things
spliced at an arbitrary joint; with it, one process runs from
abiogenesis to the present under a single conditioning, and the
technological present is a \emph{position} along it rather than a change
of subject.

The cost is that the principle is a philosophical premise and not a
result, and its scope should be stated at full width rather than left to
be inferred, because it is larger than the conditioning question that
motivated it. Continuity licenses treating a Mesozoic mammal's evasion
of predation, a body's suppression of its own cancers, and an
institution's not turning its capacity against itself as values of
\emph{one variable}. That is what permits the broader Golden program, set aside here's institutional
machinery to bear on Section~\ref{sec:imports}'s hazard dynamics at all: without
it, the adversivity model would describe the politics of societies and
the survival model would describe the fate of lineages, and no argument
would carry between them.

So the premise is not merely buying a conditioning target. \emph{It is
what makes this one work rather than several.} Should it fail, the
sections do not degrade gracefully --- they separate, and each retains
whatever it had independently while the connections lapse. \emph{Falsifier:} a demonstration that deliberative and evolutionary
hazard-comparisonn differ in a respect the conditioning must track ---
that some feature of the record is explicable under one and not the
other. The agency objection is the natural candidate and is answered
above; a better one would consist in exhibiting such a feature rather
than in restating the intuition. A second and more targeted falsifier
follows from the width just claimed: exhibit a hazard-lowering mechanism
at one scale whose governing variable provably fails to correspond to
anything at another --- a property of institutional self-restraint with
no analog in organismal or lineage-level hazard control, or the
reverse.
\end{remark}

\paragraph{The principle transfers form, and not only vocabulary.}\label{gc:rem:continuitydeductive}
One further thing the principle does is easily missed, because the
statement above is defensive in shape --- it says what continuity
licenses one to treat as a single variable, which sounds like permission
to use a word. It also licenses \emph{inference}, and an instance is now
available.

the schedule of \cite{TechnoC} represents pre-onset hazard as a sequence of
gates. Continuity says the pre-onset struggle and the post-onset one are
the same activity under different mediation. It therefore predicts that
post-onset internal hazard is representable the same way --- that the
trough channel of the survivorship companion \cite{SurvA} should decompose into gates, its
smoothness being a coarse-graining rather than a fact.
the technoanatomy companion \cite{TechnoC} carries out the decomposition and
finds two things the prediction did not have to deliver: the exponential
form falls out of geometric compounding rather than being assumed, and
the anatomy trichotomy turns out to be the schedule's own criticality
test with capacity substituted for gate count. The construction also
agrees with the boundary-scaling derivation of
the technoanatomy companion \cite{TechnoC}, which was arrived at independently.

That is the principle doing deductive work. It did not merely permit the
two hazards to be discussed in one vocabulary; it said what form one
should take given the other, and the form checked out against a
derivation that did not consume it. A reader who regards continuity as a
philosophical convenience should notice that a convenience does not
generate a correct structural prediction about a channel it was not
formulated to describe.

The falsifier stated above is unchanged and now has a sharper instance
available: exhibit an internal hazard whose dependence on capacity
provably admits no gate decomposition, and the transfer fails at exactly
the place this passage claims it succeeds.

\subsubsection{Why continuity is wanted: the non-privilege of humans}
\label{gc:subsub:nonprivilege}

The principle has now been stated and defended, but not motivated, and
its motivation is not tidiness. Continuity is wanted because it is what
makes the framework non-anthropocentric, and the non-anthropocentrism is essential rather than incidental\footnote{Provenance, in the author's words: ``The reason I decided in 2021
to explicitly build the Optimal Earth program off of conditioning on the
Golden Point/asymptotic future survival as opposed to on `observer
existence' is because `observers' felt profoundly biased as a first
principle, whereas something like `asymptotic survival capacity' felt
like something which is perfectly general enough to be fundamental on
account of not privileging humans over whatever counterfactual alien
species triggers the Golden Point in our place. The intuition behind why
this works was itself the space travel argument, where our mere apparent
contingency of being able to engage in space travel amidst a biosphere
which cannot otherwise do so seems to lock us in as an indicator toward
asymptotic survival foremost, and a present day `civilization' less
so.''}. Begin from what the target is. \emph{The Golden Point is a property of a
biosphere} --- not of a civilization, not of a species. That is why it
applies to an alien biosphere as readily as to ours, and applies without
translation: nothing in ``survival conditioning of its type'' mentions
humans, technology, or society, so nothing must be reinterpreted when
the biosphere in question is not ours.

\begin{proposition}[Recursive anchoring; \textbf{[S]}]
\label{gc:prop:recursive}
Because the target is biosphere-level and the principle is continuous
across mediation, the conditioning may be anchored at \emph{any} Aethus
in the lineage consistent with ours. Take a Mesozoic mammal, its Aethus
a parent of ours rather than an Imperfection Magnitude Drop, and ask
what the continuation of Aethic attributes within \emph{its} future
light cone must look like for its biosphere to reach the Golden Point.
The answer runs through the remainder of the lineage, and so through
us, without us having been named anywhere in the question.
\end{proposition}

\noindent Our position is therefore an \emph{output} of the conditioning
rather than a stipulated anchor for it. Humans are not a privileged
split point; the technological present is where the recursion happens to
pass, and its role is emergent from a primitive that mentions no
species. The same question asked of a mammal in the Maastrichtian, of a
lineage that has not yet speciated, or of a biosphere elsewhere entirely,
is one question.

\paragraph{Three things a human-centered target would cost.}\label{gc:rem:privilegecost}
The alternative, conditioning on ``technological civilization'' or on
observers, is not merely less elegant. It forfeits three things, in
increasing order of seriousness.

\emph{Robustness across biospheres.} An alien case would require
separate treatment, since the target would be defined by features of our
own history. A framework whose central conditioning cannot be applied to
the systems it is meant to reason about has not achieved generality; it
has assumed its instance. \emph{Robustness across perturbations.} \S\ref{gc:subsub:pastalign-argument}
evaluates histories in which humans do not arise. Under a human-centered
target those histories cannot be evaluated at all, the target simply
fails to apply, and precisely the perturbations that carry the
argument become unaskable.

\emph{Non-circularity.} A human-centered target takes \emph{us} as
primitive, so the framework would build its base unit out of the thing
it is supposed to explain. The Continuity Principle is what prevents
this, and that is the deepest reason to want it.

\paragraph{A distinction to keep: definition versus conditioning.}\label{gc:rem:defvscond}
Two claims are easily run together and must not be, since the argument's
validity depends on their separation.

The \emph{definition} of the Golden Point, a survival property of a
biosphere, in the sense Section~\ref{sec:imports} makes precise, is
stipulative.\footnote{A dated public record fixes the concept's
provenance. The author's grade-school idea journal
\cite{BenGradeSchool} records the coinage in mid-January 2021, in the
eleventh grade: ``the point after which our society will be too
capable in self preservation to regress past that point,'' with the
surrounding reasoning already survivorship-shaped --- constant danger
decaying survival exponentially, hence extinction eventually certain
without the threshold, which is the survivorship companion \cite{SurvA}(ii) in
embryo. The same record carries three further anticipations worth a
sentence each. ``Important can be defined as that which tends to
it'' is the Score of the decision-theoretic material of the broader program, set aside here as a germ. An August 2021 entry
defines ``Golden Achievement'' as the natural logarithm of the
likelihood of reaching the Golden Point, which is the additive
log-likelihood structure of the credence method
(\S\ref{gc:sub:credencemethod}) before any of its machinery existed.
And a July 2021 notation distinguishes the Miracles elapsed by the
modern day from the total ``if there are more Miracles needed in the
future'' --- the live-third-term question of
\cite{TechnoC}, posed as a subscript four years
before it was taken up as a hypothesis with an evidence route. The
record shows the conception arising from
the survivorship reasoning alone, with no contact with the
future-conditioned cosmological tradition placed at
Remark~\ref{rw:rem:futurecond}. A second dated document, the
author's \emph{Golden Science} of August 2022 \cite{GoldenScience22},
is the link between that coinage and the present formalism, and
several of this paper's core objects appear in it by name. Its
``Outlook Probability'' and ``Golden Probability'' are the
unconditioned and conditioned laws, the $\PP$/$Q$ pair, defined and
used; its internal/external division of the Departure Point is the
two-channel decomposition of Section~\ref{sec:imports}, with the external
channel argued roughly constant and the internal one fed by the
biosphere itself; its survival analysis runs on the integrated
hazard, computing a positive immortality probability $e^{a/b}$ for
exponentially decaying danger, which is the $\Lambda_\infty<\infty$
criterion in closed form; its Third Law is the inequality
$b/a<p/(1-p)$ with the L-threshold named --- the vocabulary
Proposition~\ref{gc:prop:flukes} still uses; its ``Doomsday
Pacing'' is the ancestor of the invested-timing role claim of
the occupancy discipline of the broader program; Optimation and the task-object
both appear as sections; and its ``Trace Connection,'' a
placeholder named for ``some unknown form of causal (or causal
mimicking) connection'' between our position and the Golden Point,
``until it is officially derived,'' is the promissory note that
retrocausal accordance (the retrocausal-accordance material set aside for a further paper) later paid. The same
document also contains reasoning the mature discipline now fences,
the every-attribute optimality conjectures among it, so the record
documents the corrections as well as the anticipations, which is
what a provenance record is for.} Someone who rejects everything else in this section may
adopt it, and the hazard analysis of Section~\ref{sec:imports} stands under it
regardless.

The \emph{conditioning claim}, that this property is what the
Accordance Principle has been conditioning on, is not stipulative and
is what \S\S\ref{gc:sub:conditioning}--\ref{gc:sub:humanhistory} argue
for, by model comparison on natural history and again on human history. Running
the two together would make the argument circular, defining the
conditioning target into its role and then observing that it occupies
it. Keeping them apart is what leaves the argument something it could
have failed at.

\begin{remark}[Two non-privilege claims, and their symmetry]
\label{gc:rem:twoprivilege}
The section now makes two claims of this shape, and they are the same
claim along different axes.
Principle~\ref{gc:prin:pastalign} indexes its light-cone restriction to
one's own present \emph{because that is where the observable cone is
anchored, not because the event is distinguished}; the construction at
any other event returns that event's cone.
Proposition~\ref{gc:prop:recursive} says the same of position in a
lineage: the conditioning may be anchored anywhere consistent, and
returns that anchor's continuation. No privileged spacetime event, no
privileged lineage point --- and in both cases the appearance of
privilege is an artifact of which cone the reasoner happens to occupy.
\end{remark}

\begin{remark}[A third axis: no privileged tense]\label{gc:rem:notense}
The two axes above concern \emph{where} a reasoner stands. A third
concerns \emph{when}, and it follows from the same two results rather
than needing anything new. We read the ancestral record as a settled chain in which, at juncture
after juncture, the beneficial thing happened --- which is what
conditioning predicts of any cone read from its own end. By
Proposition~\ref{gc:prop:recursive} the construction is available at any
anchor, so an observer anchored far enough downstream reads \emph{our}
actions the same way: a settled chain, cropped to their past cone, and by
Principle~\ref{gc:prin:pastalign} aligned with their own Nexus of
Optimation. We are, to that anchor, what the Cretaceous is to us. The
distinction between the static record one reads and the dynamic outcomes
one produces is therefore not a distinction in the objects. It is an
artifact of which cone the reasoner occupies --- the sentence the
previous remark ends on, now applied to tense.

Two guards, both standing elsewhere in this work and both binding here. \emph{Ensemble, not trajectory.} Nothing in this reweights anyone's
actions. Conditioning reweights the ensemble a reasoner is drawn from,
and the downstream anchor's reading says which ensemble rather than what
any biosphere will do; the broader Golden program, set aside here is disqualifying about
the other reading. \emph{And determinacy is indexed to information, not to date.} The
downstream cone holds our actions as determinate entries because that
anchor has read them; ours does not, and nothing about the reading
reaches back. That is exactly the footing on which the companion treats
Newcomb's problem, determinacy indexed to an information state rather
than to a position in time, so that an unread outcome is a blank entry
and blanks lack tense \cite{AethusDecision}, and
Remark~\ref{gc:rem:teleology}'s refusal to let a conditioning event
\emph{act} is unamended by it.

The relation to Principle~\ref{gc:prin:continuity} is close and is not
identity, and the two are easily merged.
Continuity is a claim about the \emph{object}: the ancestral struggle and
the agency-mediated one are one activity under different mediation. This
is a claim about the \emph{reading}: the epistemic relation to that
object is the same from any anchor. Neither entails the other, a
single activity could be read asymmetrically across anchors, and
readings could coincide for activities that were not one, so they are
neighbors, and the argument here consumes only the second.
\end{remark}


\section{The conditioning hypothesis: observer target against Golden target}\label{sec:conditioning}

\subsection{What alignment with the target would mean}\label{gc:sub:pastalign}

\subsubsection{Statement}\label{gc:subsub:pastalign-statement}

\begin{principle}[Past-Alignment; \textbf{[$\Phi$]}]
\label{gc:prin:pastalign}
Construct the \emph{Golden Nexus} by conditioning the Cosmic Nexus on
the Golden Point. From it, discard every stated attribute whose
referential event lies outside one's own past light cone --- in the
stricter reading, discarding the spacelike attributes as well as the
future-directed ones. What remains is the Nexus of Optimation.

Equivalently: the past light cone of Aethic attributes we actually
observe is the same one as the past light cone through our position in
the Aethic block universe of the Golden Nexus.\footnote{More precisely
than ``our actual light cone'': owing to random contingency in centric
unfolding, our own Aethus is a proper child of the Aetherian Nexus
rather than the Aetherian Nexus entire, so the object being compared is
the past light cone of the Aetherian Nexus's own Aethic block universe.}
\end{principle}

Two features of the statement are deliberate and should not be read past. \emph{The restriction is to the past cone only.} The principle takes no
position on whether attributes in one's future light cone agree with the
near-optimal part of the Golden Nexus. That silence is not modesty but
scope: one cannot reliably infer future-cone attributes prior to
engaging in Golden reasoning, so a principle that assumed agreement
there would be assuming its own conclusion.
material of this paper's earlier drafts, historical provenance only treats the forward question separately and
does not settle it.

\emph{The spacetime point is not privileged.} The principle is indexed
to one's own present because that is where the observable cone is
anchored, not because that event has any distinguished status. The same
construction at any other event returns that event's cone.

\paragraph{What the principle claims in the companion's own terms.}\label{gc:rem:oec}
Read against the Optimal-Earth Conjecture \cite{Ben26Optimal}, the
principle sharpens what ``optimal'' is optimal \emph{for}. The
conjecture there holds our record observer-optimal on effectively any
interrogable attribute. Past-Alignment says the optimand is not
vaguely observer-production but specifically the Golden Point: one asks
for the optimal Aethic history \emph{to the Golden Point}, and then
surgically reopens every attribute whose occurrence lies in the future
cone of the indexing event. If the principle holds, cropping the Golden
Nexus to one's past light cone reproduces Optimation exactly --- and it
is the exactness, not the plausibility, that would make this a
first-principles derivation of Optimation rather than a restatement.

\subsubsection{The argument from prerequisite attributes}
\label{gc:subsub:pastalign-argument}

The argument descends from the companion's axioms, and the descent is
what makes this paper's central conclusion an implication rather than
a posit. Begin inside Nexic reasoning alone \cite{Nexus}. Its
Accordance Principle prices any Aethic perturbation that replaces our
lineage with a sister clade, dinosaurs or birds or New World monkeys,
and asks what becomes of the probability that the perturbed lineage
eventually constructs rocket vehicles. The companion's answer is that
each such perturbation satisfies a Nexic Imperfection Magnitude Drop,
the joint probability of observers and the perturbed history
collapsing by the contrast the High-Contrast Earth Hypothesis
asserts; \S\ref{gc:par:dinosauroid} is the worked instance. The
construction generalizes past the physically existent sister clades
to any \emph{nameable} variation of ourselves obtained by varying a
stated set of our attributes, and no further: the companion's
biological hyperblock cosmos fixes the class of admissible
variations, every planet whether existent, hypothetical, or fictional
being representable as an alternate history of the Earth on that
canvas, and differences outside the class are not perturbations of
us at all. The comparison presupposes a common term against which a
missing capacity registers as missing, which ordinary
object-individuation cannot supply and the task-object construction (Definition~\ref{imp:taskobject}) does: what is
compared across histories is a Task-Earth and a task-lineage
(Definition~\ref{imp:taskobject}), not an Earth and a lineage.
Across the class one quantity behaves the same way. The
\emph{disposition} toward reaching another planet, a present capacity
rather than a future outcome, drops off for every nameable variant
relative to our own value for it, whether the variant is a persisting
dinosaur lineage that never reaches the cognitive package, a fungal
world lacking the animal capacities the construction requires, or a
great-ape lineage left arboreal by a different African climatic
sequence. That is a deduction of Nexic reasoning about ourselves, and
it is the \emph{pinnacle invariant}: among the nameable variations of
ourselves, we hold the highest disposition toward reaching other
planets, because each alternative has a lower chance of evolving far
enough to acquire the capacity at all. The claim is local by
construction, ranging over variations of us and not over the space of
describable beings.

Golden reasoning takes the invariant as its deductive root, and the
move is a substitution into the companion's own principle. Two
features of the Accordance Principle as the companion states it make
the substitution available. The principle is stated at a general
conditioning target: the draw is asserted between the two middle
Aethae of a containment chain
$H^{\ast}\sqsubseteq H^{-}\sqsubseteq H^{+}\sqsubseteq H$ and not
between the record's ends, so that $H^{+}$ ranges over targets rather
than being fixed at observer-production, and the principle holds at
each. And the chain is tense-free: $\sqsubseteq$ is Aethic
containment and not temporal order, nothing in the chain requiring one
term to fall earlier or later than another in the history's own time,
which is the Indiscriminate Time Principle carried into the draw
(Remark~\ref{gc:rem:notense}). A future-referring conditioning target
is one of the targets the principle was stated to range over.

Make the substitution. Let $A$ be the ability to travel in space, a
realized attribute of our Aethus and a present disposition, with $A'$
its canonical complement; and let $H^{+}$ be the ability to
systematically counter cosmic threats, which is the nullification of
the external Departure hazard of Section~\ref{sec:imports} and so is the Golden
Point, of which travel in space is a necessary requisite. The
Accordance Principle then reads, with no step added,
\begin{equation}\label{gc:eq:pinnacle}
\frac{\PP(\text{Golden Point}\mid A')}{\PP(\text{Golden Point}\mid A)}
\;\lesssim\;\frac{\PP(A)}{1-\PP(A)},
\end{equation}
in the exceedance modality the companion's derivation carries, a
$\kappa$-fold shortfall confined to probability $1/(1+\kappa)$. The
right-hand side is where the Nexic deduction enters as the empirical
input: $\PP(A)$, the weight of the space-reaching disposition under
the Cosmic Nexus, is tiny, because across the nameable variations of
ourselves the disposition drops off. So futures that reach the Golden
Point without our space-reaching capacity are disfavored, relative to
futures that reach it with that capacity, by that ratio. That
inequality is the Past-Alignment principle at one attribute
\emph{under the Golden-conditioned draw}: it states what the record
looks like if it is drawn under conditioning on the target, and it is
deduced from the companion's principle with that outer measure
supplied. Whether our record is so drawn is the hypothesis $H_G$, and
the evidence for it is between-model, the Bayes factor of
\eqref{gc:eq:bfGO}, with the prior bound's tail as the within-model
diagnostic; in the critical regime the target is read as the horizon
family $\{T>t\}$, per Theorem~\ref{imp:targeted}.

The substitution runs at any Golden-relevant trait, not only at space
travel: any trait of ours may be placed in the $A$ seat, and
\eqref{gc:eq:pinnacle} then prices its complement at the trait's own
Cosmic-Nexus weight, provided the trait is identifiable in histories
where it does not occur, which is the task-object identification of
Definition~\ref{imp:taskobject}. Summed over the traits so interrogated, the
result is a local statement about the record.

\begin{corollary}[Local Golden alignment; \textbf{[S]}]
\label{gc:cor:localopt}
If for every Golden-relevant trait, individuated as a task-object per
the task-object construction (Definition~\ref{imp:taskobject}), the perturbed histories depress its Cosmic-Nexus
weight, then by \eqref{gc:eq:pinnacle} applied at each trait every
tested neighboring perturbation of our Aethus is Golden-disfavored:
our Aethus is a coordinatewise local maximizer of the Golden Point's
probability among its interrogated neighbors. \emph{Strong
Past-Alignment}, that the whole observed past cone is the restriction
of a globally Golden-optimal history, does not follow: a landscape can
be worse under every one-coordinate move and better under some
multi-coordinate one. Strong Past-Alignment is conditional on the
cumulative-suppression landscape assumption the Optimal-Earth companion
makes explicit \cite{Ben26Optimal}, and is carried at that assumption's
grade, not at this corollary's.
\end{corollary}

\begin{remark}[The grade]\label{gc:rem:grade}
The corollary is \textbf{[S]} and not \textbf{[P]}, for two reasons
that are separable. Its inference form is the companion's and carries
the companion's grade; what it adds is the chain condition, that a
reasoner's Aethus stands in the containment relation to a
future-referring target, which is the Route World's own commitment
and is graded where that commitment is made
(\cite{TargetedB1}). And it establishes local optimality
among interrogated perturbations, not global optimality over the
space of histories. What would strengthen it is a perturbation
exhibited that leaves a Golden-relevant trait's weight undiminished;
what would break it is a systematic class of such perturbations.
\end{remark}

Three things follow about the architecture. The Retrocausal Accordance
Principle of \cite{TargetedB1} is the companion's principle
at a tense-free target, its display Theorem~\ref{imp:targeted} being
\eqref{gc:eq:pinnacle} with the target left general; its
\textbf{[$\Phi$]} grade carries the chain condition and not the
inference form. The non-circularity is structural: the tense-freeness
that licenses the substitution is the companion's, inherited from the
Aethic Indiscriminate Time Principle, so Golden reasoning borrows no
conclusion of its own, and the rarity $\PP(A)$ that fills the
right-hand side is deduced under the Cosmic Nexus with no reference to
any future. And the choice of target is not free. One might abstract
from the space-reaching invariant to maximal capability or to maximal
information processing and obtain a different $H^{+}$;
Departure-nullification is selected because the survivorship
mathematics of Part~II independently picks it out as the fixed point
of \S\ref{gc:subsub:cascade}, the target at which the external channel
is driven to zero.

A second consideration selects the same target. The nearby candidate
is ``technologically advanced civilization,'' and it fails twice: it
admits no principled threshold, and it names a human societal
construct with no continuous extension backward through the record
the argument must run on. The Golden Point is primitive enough to
avoid both failures, being survival conditioning of a stated type, and
so stands in substantive continuity with natural history: the same
target that a Mesozoic mammal's persistence bears on bears on a launch
vehicle, which is Principle~\ref{gc:prin:continuity} and the reason
the task-object individuation of the task-object construction (Definition~\ref{imp:taskobject}) is available at every
rung. A target one cannot define across the record cannot be
conditioned on across the record.

The paper's fundamental conclusion is a further implication of this
move and not a separate posit. That our past record is effectively
conditioned on our asymptotic-future survival, the Past-Alignment
principle of \S\ref{gc:sub:conditioning}, the retrocausal instrument
of Theorem~\ref{imp:targeted}, and the strong reading of the
Route World all descend from the pinnacle: if the record is such that
every nameable variation of us is worse placed to reach other planets,
and reaching other planets is a requisite of the target, then the
record reads as conditioned on the target. The pinnacle invariant is
the inferential root of the Past-Alignment principle.


\subsection{The conditioning principle, and the two arguments for it}\label{gc:sub:conditioning}

Past-Alignment as stated is an observation about agreement. The
following argues that the agreement is not coincidental --- that
Nexic natural history was never conditioned on Optimation in the first
place, but on the Golden Nexus throughout.

\begin{principle}[Golden-Conditioning; \textbf{[S]}, by model comparison]
\label{gc:prin:conditioning}
Two hypotheses about the outer measure of the record are compared:
$H_O$, that the record is drawn under conditioning on personal
observer-existence, and $H_G$, that it is drawn under conditioning on
the Golden Point. A low-probability realization under $H_O$ is not a
contradiction of $H_O$; it is evidence, and the argument is stated as
evidence. Consider first what $H_O$ would predict.

Observerhood is robust under Aethic perturbation. One may strip
probabilistic prerequisites of the Golden Point, rocket capability,
the presence of industrialization-driving coal, and replace them by
their complements without disturbing one's standing as an observer. A
world in which Earth simply lacked the resources to begin an industrial
revolution is a world whose inhabitants are no less observers. Now select a handful of such attributes: necessary for the Golden Point,
unlikely to survive reopening in the Aetherian Nexus, and admissible in
the sense of Remark~\ref{gc:rem:boundrequires} --- reopening them must
leave the dependent structure intact, so that the resulting Aethae
remain individuated by the same task-objects as ours
(Definition~\ref{imp:taskobject}). The totality of
the counterfactual Aethae so obtained carries \emph{substantially more}
Aethic weight than the Aetherian Nexus, while carrying a comparatively
minute probability of the Golden Point.

Under $H_O$ the record is a draw from that high-weight counterfactual
sum, and ours is a low-weight member of it; under $H_G$ the record is a
draw from the members carrying the target's prerequisites, of which ours
is typical. The conclusion is a Bayes factor, and nothing stronger. Let $E$ be the
conjunction of the technological-route prerequisites so selected, with
$\PP(E\mid H_O)=q\ll1$ because observerhood is insensitive to them and
$\PP(E\mid H_G)\approx1$ because the target requires them. Then
\begin{equation}\label{gc:eq:bfGO}
\mathrm{BF}_{G:O}(E)=\frac{\PP(E\mid H_G)}{\PP(E\mid H_O)}\approx\frac1q,
\end{equation}
which is enormous when $q$ is small and is exactly how much evidence
the record supplies for the Golden target over the observer target.
The principle asserts that this factor favors $H_G$ by orders of
magnitude, and the weight bound of \S\ref{gc:sub:naturalhistory} is
the estimate of $q$ that gives the factor its size. Both hypotheses are
compared after conditioning on the common baseline $C$ the record
already fixes, that observers exist now and that the biosphere has
survived to the present, so that $\mathrm{BF}_{G:O}(E)=\PP(E\mid
C,H_G)/\PP(E\mid C,H_O)$ and no part of the factor is earned by
evidence both hypotheses already had to explain.
\end{principle}

\subsection{The undeniable case: human history}\label{gc:sub:humanhistory}

Principle~\ref{gc:prin:conditioning} admits two arguments, and they
differ in one respect that decides their order. Both need observerhood
to be robust under the perturbations considered --- that one can strip
prerequisites of the Golden Point without disturbing anyone's standing
as an observer. On natural history a sufficiently demanding notion of
``observer'' can contest that premise (Remark~\ref{gc:rem:circularity}).
On human history it cannot: no one holds that human observers required
Alexander's survival, or Newcastle's coal, or Newton. Observers existed
before each and would have continued after. So the human-history
argument is the one on which the robustness premise is simply not
deniable, and it is put first for that reason; the natural-history
argument follows as the same reasoning in a domain where the premise is
contestable and the conclusion correspondingly weaker.

The argument also does something the natural-history one cannot. What
these contingencies deliver is not observers but specifically the
technological condition --- and observer-conditioning is indifferent to
every one of them while Golden-conditioning is not. That isolates what
the conditioning must be sensitive to, and it says something about the
companion that the companion cannot say about itself: its
conditioning target, an eventual observer in one's biosphere, is exactly
right for the pre-technological record and provably insufficient beyond
it. Not a criticism, the companion reads the pre-human record and
stops, by design, but the reason for stopping becomes visible only
here.

\subsubsection{Flukes establish their own significance}
\label{gc:subsub:flukes}

\begin{proposition}[Significance from route-deletion; \textbf{[S]}]
\label{gc:prop:flukes}
Let $F$ be an event established as a \emph{fluke}, of Cosmic-Nexus
weight $p<\tfrac12$, whose complement $\neg F$ is route-deleting in the
sense of Remark~\ref{gc:rem:flukescope}. Under the Golden-conditioned
draw, the joint event that $F$ is realized and is the less
Golden-productive member of its pair has probability at most $p$
\eqref{gc:eq:priorbound}: a rare realized route-deleting attribute is
unlikely also to be Golden-unproductive. The statement is within-model.
Establishing that $\neg F$ is route-deleting rather than a permutation
is the case work of Remark~\ref{gc:rem:flukescope}, and evidence that
the record is drawn under the Golden target at all is the between-model
factor \eqref{gc:eq:bfGO}.
\end{proposition}

\noindent The proposition's core is the companion's principle at one
choice of $\kappa$, and it carries neither the Second Law's
translation nor a separate self-sampling step. Taking
$\kappa=(1-p)/p$ in the Accordance display collapses its right-hand
side to unity and its bound to $p$, so that
\begin{equation}\label{gc:eq:priorbound}
\Pr\!\big[\,\PP(G\mid\neg F)\;\ge\;\PP(G\mid F)\,\big]\;\le\;p:
\end{equation}
under the Golden-conditioned draw, the joint event that the rare member
is realized and is the less Golden-productive of its pair has
probability at most $p$: fix the member $A$ of prior $p<\tfrac12$;
$\PP(G\mid A')\ge\PP(G\mid A)$ holds exactly when the $G$-conditioned
weight $w$ of $A$ satisfies $w\le p$, and under $\PP_G$ the member $A$
is realized with probability $w$, so the joint event has probability
$w\,\one\{w\le p\}\le p$. That is the claim that flukes establish
their own significance, and it is a \emph{within-model} diagnostic: it
says what a rare realized attribute is, given that the record is drawn
under conditioning on the target. It is not \emph{between-model}
evidence that the record is so drawn; that evidence is the Bayes factor
of \eqref{gc:eq:bfGO} and Proposition~\ref{gc:prop:credence}, and the
two must not be run together. What the case work of
Remark~\ref{gc:rem:flukescope} supplies is not the significance but
the \emph{pairing}: that $\neg F$ is the canonical complement rather
than a permutation, so that the attribute whose prior is $p$ is the
one whose significance \eqref{gc:eq:priorbound} prices.

\begin{proposition}[The stockpile: accordance on the conjunction;
\textbf{[S]}]\label{gc:prop:stockpile}
Let $F_1,\dots,F_n$ be a \emph{predeclared} family of route-deleting
flukes with Cosmic-Nexus weights $p_1,\dots,p_n$, and let
$F=\bigcap_iF_i$, with complement $\neg F=\bigcup_i\neg F_i$, the union
of the counterfactual Aethae. Suppose the family is submultiplicative,
$\PP(F)\le\prod_ip_i$, which holds with equality under independence and
holds under negative association; positive association gives the reverse
inequality and defeats the bound. Then under the Golden-conditioned draw
\begin{equation}\label{gc:eq:stockpile}
\PP_G\bigl[\,F\text{ realized and }\PP(G\mid\neg F)\ge\PP(G\mid F)\,\bigr]
\;\le\;\PP(F)\;\le\;\prod_ip_i,
\end{equation}
the joint-event form of \eqref{gc:eq:priorbound} applied to the single
attribute $F$: the bound tightens multiplicatively in $n$ while the
counterfactual side accumulates weight $1-\prod_ip_i$, so that ten
flukes at $p=0.1$ each price the joint event at $10^{-10}$ against a
counterfactual union of weight $1-10^{-10}$.
\end{proposition}

\noindent The aggregation is where the argument's force lives, and two
conditions govern it. Each member must pass the
route-deletion audit of Remark~\ref{gc:rem:flukescope} separately,
since a permutation admitted into the conjunction inflates the bound
without adding linkage. And the members must be interrogated under a
declared scan count, the multiplicity guard of
the multiplicity guard of \cite{TargetedB1} applying to the stockpile as to any member.
The guard has two forms and they are kept apart: the expected number of
spurious $\kappa$-fold flags among $n$ interrogated pairs is at most
$n/(1+\kappa)$, and the tail is bounded without any independence
assumption by Markov, $\Pr[N\ge m]\le n/(m(1+\kappa))$, with the
familywise form $\Pr[N\ge1]\le n/(1+\kappa)$ by the union bound.
Exceeding the expectation is not significance; a declared family is
scored against the tail.

\begin{remark}[What the proposition quantifies over, and the
counterfactuals it excludes]\label{gc:rem:flukescope}
The template above, $F$ is improbable, we observe it, therefore $F$
is linked to the Golden Point, runs on \emph{any} observed
improbability if nothing restricts the range of $F$, and would then
license absurdities: the precise configuration of a Cretaceous
raindrop is astronomically improbable and is not thereby significant.
The restriction that blocks this governs the proposition's use and is
stated before the stockpile that depends on it.

\emph{The counterfactual must delete a route, not permute one.} Write
$\neg F$ for the alternative. The inference is licensed when
$\PP(G\mid\neg F)$ is genuinely lowered by the substitution, when
the counterfactual removes a prerequisite of the Golden Point, and
is void when $\neg F$ merely reshuffles which particular events fill
the same structural role. The raindrop case is a permutation: some
configuration obtains, and every one of them leaves the downstream
dependency structure untouched, so $\PP(G\mid\neg F)=\PP(G\mid F)$ and
the bound material of this paper's earlier drafts, historical provenance only is satisfied with room to spare
rather than violated. Improbability alone therefore establishes
nothing; improbability \emph{of a route-deleting attribute}
establishes the linkage, and the proposition should be read with that
restriction throughout.

\emph{The restriction is not free, and it is what the stockpile is
selected for.} Deciding whether a counterfactual deletes or permutes
is a claim about dependency structure and is not delivered by the
attribute's weight --- which is
Remark~\ref{gc:rem:flattening}'s point arriving early: a marginal test
on weights recovers the hubs of the dependency graph, and the
route-deleting attributes are exactly the hubs. So the criterion is
principled and its application is \textbf{[S]}, case by case. This is
also the honest register in which the stockpile's items differ. The
coal has the route-deleting character on the economic historians' own
analysis, Allen's price ratios make mechanization pay in one place
and nowhere else, so the counterfactual removes the route rather than
relocating it \cite{Allen09}, while the Granicus is weaker on
exactly this axis, since the claim that no comparable transmission
survives Alexander's death at the Granicus is a strong fragility
premise about the Hellenistic world and not a computed dependency.
That ordering is the same one Remark~\ref{gc:rem:narrativebias} reaches
from the evidential side, and the agreement of two independent
criteria on which items are strongest is worth more than either.
\end{remark}

\begin{remark}[The retrospective-target null, dissolved by the principle]
\label{gc:rem:sharpshooter}
The strongest rival to the proposition is not the determinist's
substitutability but the \emph{retrospective-target} null: any
realized history, read backward, contains improbable events that were
prerequisites of whatever that history's distinctive features happen
to be, so a stockpile of flukes pointing at the realized outcome is
expected under no conditioning whatever, the sharpshooter who paints
the target around the shots. The null is answered not by argument but
by the form of the Retrocausal Accordance Principle
(Theorem~\ref{imp:targeted}), which was built so that the
two freedoms the fallacy exploits are closed and its residual content
is priced. The first freedom is the choice of target, and the
principle's $B$ is not drawn from the record: it is the Golden Point,
picked out by the fixed-point structure of \S\ref{gc:subsub:cascade}
as a future-cone attribute defined before and independently of which
record obtains, so the target stands before the shots in the only
sense the null requires. The second freedom is the choice of
reference class, and the principle's $A'$ is the canonical complement
of the realized attribute, so nothing about the alternatives is
analyst-supplied and no class can be drawn to make the realized value
look rare. What the null correctly observes is then exactly what the
principle's exceedance form Theorem~\ref{imp:targeted} quantifies:
rare attributes that are irrelevant to $B$ do abound, they violate the
bound, and the principle prices each such violation at probability at
most $1/(1+\kappa)$, so that among $n$ interrogated attributes up to
$n/(1+\kappa)$ spurious $\kappa$-fold flags are expected, which is the
noise floor of the multiplicity guard of \cite{TargetedB1}. The null is thereby
absorbed rather than refuted: its content is a term in the
principle's own accounting, and a trace standing clear of the floor
is evidence while one within it is the expected noise.

Three further answers to the null are corollaries of the same
form. That consistency with a hypothesis is not evidence for it
over a rival the same data fit, the comparative point of
\S\ref{gc:sub:historicalcases}, is the principle's left side read
literally: an attribute is evidence only through the ratio
$\PP(B\mid A')/\PP(B\mid A)$, and an attribute whose complement leaves
$B$ as likely as before fails the bound and is counted as noise,
which is the route-deletion-against-permutation question of
Remark~\ref{gc:rem:flukescope} stated as arithmetic. And that the
human-history leg's distinctive content lies in flukes that are
prerequisites of the technological condition specifically, to which
observer-conditioning is provably indifferent, is the identification
of which attributes have $\PP(B\mid A')\ll\PP(B\mid A)$ under the
Golden target and not under the observer target: it is where the
ratio discriminates the two conditionings, and raw improbability
never was the currency. One discipline the dissolution does not
remove, because it is what the floor is for: the scan count must be
declared. The null returns, in full force, against any reasoner who
interrogates attributes freely and reports only the flags, and the
principle's guard is an instrument only when $n$ is on the table.
\end{remark}

\paragraph{The dinosauroid, answered (worked case; \textbf{[S]}, one
import at \textbf{[$\Phi$]}).}\label{gc:par:dinosauroid}
The question whether the dinosaurs could have evolved into humans,
posed in its classic form as the dinosauroid \cite{RussellSeguin82}
and generalized in the debate between contingency and convergence
\cite{Gould89,ConwayMorris03}, receives a direct answer once its
vague target is replaced. ``Evolving to intelligence'' is not a
well-posed endpoint; ``acquiring the capacity to trigger the Golden
Point'' is, and with that substitution the framework's answer is no,
by deduction rather than by paleontological intuition. The classic
rejoinder, that some birds are strikingly intelligent, is met on its
own terms: a lineage may become smart along one metric or another
without ever rising to high capacity for high Golden-Point
probability, and the delicate Golden-leading activities, rocketry as
the emblem, would never have entered its realm of possibility. The
companion's Imperfection-Magnitude-Drop argument \cite{Nexus}, an
IMD being a targeted perturbation of one's own natural history whose
observer-production probability is drastically minute, sharpens
``less'' to ``vastly less'': such a lineage's Golden-Point probability would sit far
below our own, so that with the vast majority of the likelihood it
would have become smart, if at all, along a tangential metric. The
substitution also dissolves the standoff it was borrowed from.
Convergence may well hold for cognition along some metric, as its
defenders argue, while contingency holds at the level that matters
here, because Golden capacity is a delicate, Miracle-coupled
competence and not a convergent target; the two camps were arguing
about different endpoints under one word.

The case is the flukes proposition's canonical instance, and it
shows the route-deletion criterion doing its case work rather than
being asserted; it is also the worked instance of the pinnacle
invariant from which \S\ref{gc:subsub:pastalign-argument} descends
the conditioning principle. The end-Cretaceous impact is the fluke $F$, the
world in which the dinosaurs continue is $\neg F$, and the claim
that $\neg F$ carries vastly lower Golden-Point probability is the
claim that $F$ was route-deleting rather than a permutation, which
is exactly what Remark~\ref{gc:rem:flukescope} requires to be
established case by case. The ``birds'' objection is the
substitutability reading of Remark~\ref{gc:rem:determinism},
answered here by the metric distinction. One pitfall must be
avoided, and the companion names it: the naive counterfactual ``no
humans'' removes the conditioning on us, and with it Optimation, so
that the counterfactual Earth destabilizes quickly and the
comparison becomes trivial. The comparison is therefore run already
conditioning on Earth maintaining its habitability, and the answer
survives that concession, which is what gives it force. The
structure is the natural-historical mirror of the extinction
proposition (the broader Golden program, set aside here): remove the muffling lineage
and the background stands unmuffled, in the deep past as in the
counterfactual present.


\begin{remark}[This settles the contingency-determinism debate, in one
direction]\label{gc:rem:determinism}
Historians divide over whether the path to modernity was contingent or
overdetermined, and the dispute is usually conducted by weighing
narratives. Proposition~\ref{gc:prop:flukes} makes it an inference. The determinist holds that the apparent flukes were substitutable: no
Alexander, then someone else; no British coal, then some other route.
Substitutability is exactly high $b$, a good chance of the Golden
Point without $F$, hence low Golden Linkage, hence $g \approx p$,
hence the expectation that we occupy $\neg F$. The Third Law, the
Accordance bound material of this paper's earlier drafts, historical provenance only with $H$ the Golden Point, bounds it quantitatively: $b/a < p/(1-p)$, so the flukier the event, the more
severely substitutability is constrained by our observing it. The
determinist is not refuted by a better narrative; they are constrained
by an inequality whose inputs are the event's improbability and its
route-deleting character (Remark~\ref{gc:rem:flukescope}) --- nothing
else.
\end{remark}

\subsubsection{The stockpile}\label{gc:subsub:stockpile}

Two cases, deliberately ordered so that the weaker rhetorically comes
first. \emph{A single sword-stroke.} At the Granicus in 334 BCE, two months
into an eleven-year campaign, the Persian noble Spithridates brought a
battle-axe down on Alexander's helmet and split its crest; Cleitus the
Black severed his arm before the second blow \cite{Arrian,
PlutarchAlex}. Had the stroke landed, the expedition would have ended
within days of beginning, and with it the Hellenistic world, its
Alexandrias, and the transmission through which much of what followed
ran.

\emph{Where the coal happened to be.} Pomeranz argues that England was
a ``fortunate freak'' whose coal lay close to water and ports, making
steam economically feasible, while China's coal sat in the northwest,
far from the Jiangnan textile economy that might have used it ---
a geological contingency placing coal and the Americas nearer the
western than the eastern end of Eurasia \cite{Pomeranz00}. Allen adds
the labor side: Newcastle carried the highest ratio of labor to energy
costs in the world in the early eighteenth century, and only under that
ratio did mechanization pay \cite{Allen09}. His summary judgment is
an Imperfection Magnitude Drop stated by an economic historian who has
never encountered one: there was a single path to the twentieth century,
and it ran through northern Britain.

A third case is available but requires reframing. Diamond's account of
Eurasian advantage \cite{Diamond97} is \emph{determinist} at the level
of human history, geography made the outcome near-inevitable, and
is criticized on that ground. What the present argument takes from it
lies one level down: the distribution of domesticable species and the
continent's east--west axis are contingent \emph{across possible
Earths} while determining \emph{within} ours. That is precisely the
Accordance profile, low outlook probability with high Golden Linkage,
and Diamond's determinism is at the wrong level rather than mistaken.

\subsubsection{Why this leg is the stronger one}
\label{gc:subsub:strongerleg}

Two reasons, and the second is the more important. \emph{The robustness premise becomes uncontroversial.} The natural-history
comparison needs it granted that observers survive the removal of the
end-Cretaceous impact, which can be resisted. Nothing comparable is
available here. No one holds that human observers required Alexander's
survival, or Newcastle's coal seams, or Newton. Humans existed for
millennia before each and would have continued after their removal. The
premise that carried the weight and could be denied is here simply not
deniable.

\emph{The conclusion sharpens.} What these contingencies deliver is not
observers but specifically the \emph{technological} condition ---
skyscrapers, the internet, the capacity to leave a dying star.
Observer-conditioning is indifferent to every one of them;
Golden-conditioning is not. So the human-history argument does not merely repeat
the first at lower risk; it isolates what the conditioning must be
sensitive to.

\begin{remark}[The narrative-bias objection, and the test that answers
it]\label{gc:rem:narrativebias}
The serious objection is selection in the record rather than in the
world. Dramatic near-misses make better stories, are more often written
down, and are more often repeated; a history dense with them may report
a fact about historiography. The framework predicts an asymmetry that narrative bias does not, and
the prediction is already on record. Attributes with high Golden Linkage
magnitude but \emph{low connectivity} to other attributes should have
low Golden Probability --- lucky avoidance of adversity outranks lucky
survival of it, so isolated close calls should be \emph{rare} in the
record while densely connected contingencies should be
\emph{numerous} \cite{GoldenScience22}. Narrative bias inflates both
alike; the framework inflates one and suppresses the other.

The discrimination favors the framework for a plain reason: \emph{the strongest cases are the dull ones}. Nobody narrates
the Newcastle labor-to-energy price ratio, and no popular account turns
on Jiangnan's distance from northwestern coal seams. Yet that is where
the peer-reviewed economic history converges \cite{Pomeranz00,
Allen09}. Were narrative selection generating the pattern, the
undramatic structural contingencies would be missing --- and they are
the best-attested items in the file. The Granicus is the vivid
illustration; the coal is the evidence, and the ordering of
\S\ref{gc:subsub:stockpile} inverts the rhetorical one deliberately.
\end{remark}


Both legs concern the past cone, and neither speaks to the future. That
silence is deliberate, \S\ref{gc:subsub:pastalign-statement} recorded
it as scope rather than modesty, and the three positions available on
the forward question are set out next without adjudication.


\section{The natural-history leg: the counterfactual weight, bounded}\label{sec:natural}

\subsection{The contestable case: natural history}\label{gc:sub:naturalhistory}

The same comparison, run on the natural-historical record. Its premise, that observerhood survives the perturbations considered, is
contestable here in a way it was not above, and the argument is
correspondingly the weaker of the two; it is included because it reaches
further back, and because it is where the conditioning claim was first
posed.

\paragraph{The Trace Connection, named in 2022.}\label{gc:rem:trace}
The claim this subsection advances was posed as an open problem, under
its own name, four years before the present argument. Our
\emph{Golden Science} of August 2022 \cite{GoldenScience22} sets out the
Golden Principle in four numbered points, of which the third states that
our observations are far closer to representing those of an average
Golden-Point biosphere than those of an average position in the universe, Past-Alignment in prose, and the fourth states that some
connection, causal or causal-mimicking, holds between our existence at
this time and place and the Golden Point itself. That paper then names
the connection and declines to supply it: it is to be called \emph{the
Trace Connection} until it is officially derived.

\S\ref{gc:sub:conditioning} is offered as that derivation. What the 2022
statement supplies, and what no later argument could manufacture for
itself, is that the gap was identified, named, and dated before any
candidate filling existed --- so the comparison cannot be read as having
been reverse-engineered to fit a conclusion already in hand. The same
document also carries, in prose and without the stochastic apparatus of
Section~\ref{sec:imports}: the hazard-function framing of biosphere survival; the
observation that an exponentially decaying hazard yields a strictly
positive limiting survival probability; the division of the failure mode
into internal and external channels, together with the insistence that a
Golden biosphere must lessen \emph{both} rather than either; the thesis
that domination is inefficient because it feeds the internal channel,
which is $\Phi3$; the Kardashev critique including the
self-consumption argument; and the definition of the Golden Point as the
point past which a biosphere cannot regress, drawn there against the
Hubble horizon. The formal content of
Section~\ref{sec:imports}--the broader Golden program, set aside here is later; the apparatus it formalizes is
not.

\begin{corollary}[What remains; \textbf{[S]}]\label{gc:cor:remains}
What the comparison leaves standing is the empirical artifact of
Past-Alignment itself: conditioning on the existence of the Golden Point
somewhere in the Aethic block universe reproduces exactly the
inferential form of the Accordance Principle that yields the
stated-attributes of Nexic reasoning. Among the contingent natural
histories one might find oneself at the end of, a fungal world, a
dinosaurs-survive world, a wholly lemurine one, the one that must be
conditioned on to reproduce the Nexic state of affairs we actually
observe is the one in which the Golden Point exists.
\end{corollary}

\begin{remark}[A second robustness, and why the target's modal form was
the right one]\label{gc:rem:finitehorizon}
Principle~\ref{gc:prin:conditioning} turns on a robustness: observerhood
survives perturbations the record does not, so a conditioning
indifferent to the difference cannot have produced the difference. The
schedule channel of the schedule of \cite{TechnoC} exhibits a second target
with the same defect, and it is nearer to hand than observerhood. \emph{Survival to a finite horizon is robust in exactly that way.} A
biosphere may survive any stretch one names without having climbed its
schedule at all --- by not being struck, at bare background, on the
lucky branch that the technoanatomy companion \cite{TechnoC} shows never leaves the
support of a finite-horizon conditioning. So conditioning on
having-lasted leaves the record unexplained for precisely the reason
conditioning on observers did: it is satisfied by histories in which
nothing was climbed.

What survives the objection is the asymptotic form, and the point worth
extracting is that the target was already stated in it.
Definition~\ref{gc:def:departure} gives Goldenness as $\exists\forall$, some course admits \emph{no} foreclosing continuation, which is a
condition on the whole of a future and not on any initial segment of
one. A reading of the Golden Point as a passage completed at some finite
time would have inherited the robustness above and failed here. The
modal form is therefore not a technical nicety of the definition; it is
what makes the target capable of doing the work
Principle~\ref{gc:prin:conditioning} asks of it.

The cost is recorded rather than absorbed: this raises what the
conditioning must be, and by Theorem~\ref{imp:ceiling} no one is entitled
to know that it obtains.
\end{remark}

\subsubsection{Bounding the counterfactual weight}
\label{gc:subsub:weightbound}

The comparison of \ref{gc:prin:conditioning} turns on a single inequality:
that the counterfactual Aethae obtained by reopening Golden-Point
prerequisites carry substantially more weight than the Aetherian Nexus.
Stated bare, that is an assertion, and an objector may simply deny it ---
perhaps coal is common given a Carboniferous, perhaps the metallurgical
sequence is near-forced given hands and fire. This subsection replaces
the assertion with a construction, a bound, and a statement of exactly
what empirical input remains outstanding.

\paragraph{The construction.} Fix stated-attributes
$a_1,\dots,a_k$, each a prerequisite of the Golden Point and none a
prerequisite of observerhood. Let $\mathcal A$ denote the Aetherian
Nexus, in which every $a_j$ obtains, and let $\mathcal A_j$ denote the
Aethus obtained by reopening $a_j$ alone to its logical complement,
holding the remaining $k-1$ at their realized values.

\begin{lemma}[Disjointness; \textbf{[P]}]\label{gc:lem:disjoint}
The Aethae $\mathcal A_1,\dots,\mathcal A_k$ are pairwise disjoint, and
each is disjoint from $\mathcal A$. For $j\neq j'$, the pair
$\mathcal A_j$ and $\mathcal A_{j'}$ disagree on \emph{both} slot $j$
and slot $j'$; and $\mathcal A_j$ disagrees with $\mathcal A$ on slot
$j$.
\end{lemma}

\begin{proposition}[Weight ratio; \textbf{[P]}]\label{gc:prop:ratio}
Write $\alpha_j$ for the probability that slot $j$ takes its realized
value, conditional on the remaining slots taking theirs. Then
\begin{equation}\label{gc:eq:ratio}
\frac{W(\mathcal A_j)}{W(\mathcal A)}=\frac{1-\alpha_j}{\alpha_j},
\end{equation}
and by Lemma~\ref{gc:lem:disjoint} the single-flip Aethae may be summed:
\begin{equation}\label{gc:eq:bound}
R \;:=\; \frac{W_\times}{W(\mathcal A)}
\;\ge\; R_1 \;=\; \sum_{j=1}^{k}\frac{1-\alpha_j}{\alpha_j}.
\end{equation}
\end{proposition}

\begin{remark}[What the bound does and does not require]
\label{gc:rem:boundrequires}
Three features of \eqref{gc:eq:bound} are worth isolating, because they
are what make it usable.

\emph{Independence is not required.} Only disjointness is, and that
holds by construction. Dependence among the $a_j$ changes what the
individual ratios \emph{are}, the $\alpha_j$ become conditional
rather than marginal, as stated, but not whether the terms may be
added. This matters because the attributes at issue are plainly not
independent, and an argument needing them to be would be unusable. \emph{It is a lower bound, and deliberately a loose one.} The
single-flip Aethae are a vanishing fraction of the counterfactual set,
and the full sum over all Aethae differing in at least one slot is
larger by many orders; it has no closed form for dependent attributes,
since $\prod_j\PP(a_j\mid a_{-j})$ is not $\PP(a_1,\dots,a_k)$ without
a factorization assumption, and none is imposed. The comparison needs
only that $R\gg1$, which the single-flip sum supplies exactly.

\emph{Admissibility is governed by the Imperfection Magnitude Drop
criterion.} A perturbation enters the sum only if it leaves the
dependent structure intact --- if reopening $a_j$ does not re-randomise
what depends on it. A perturbation that triggers a drop does not yield a
clean single-slot flip, and Lemma~\ref{gc:lem:disjoint} fails for it.
Coal in an inaccessible location is admissible on this criterion: it may
be probabilistically catastrophic without re-randomizing the structure
downstream. A relocation of the lineage's speciation to another
continent is not, since the material that later mediates the outcome is
not present there to be held fixed.
\end{remark}

\paragraph{How large must $R$ be, and how improbable must the
attributes be?} The comparison requires only $R\gg1$. Evaluating
\eqref{gc:eq:bound} for $k$ attributes each at probability $\alpha$:

\begin{center}\small
\begin{tabular}{@{}rrrrr@{}}
\toprule
$\alpha$ & $k=1$ & $k=3$ & $k=5$ & $k=8$\\
\midrule
$10^{-1}$ & $9$ & $27$ & $45$ & $72$\\
$10^{-2}$ & $99$ & $297$ & $495$ & $792$\\
$10^{-3}$ & $999$ & $3.0\times10^{3}$ & $5.0\times10^{3}$ & $8.0\times10^{3}$\\
$10^{-4}$ & $1.0\times10^{4}$ & $3.0\times10^{4}$ & $5.0\times10^{4}$ & $8.0\times10^{4}$\\
\bottomrule
\end{tabular}
\end{center}

\noindent The bar is far lower than the informal argument suggested. A
\emph{single} attribute at $\alpha=10^{-3}$ delivers $R\sim10^{3}$
unaided; no portfolio is needed, and the burden reduces to establishing
one prerequisite as genuinely improbable. Correspondingly, the objection
that the counterfactual sum might be comparable requires \emph{every}
candidate prerequisite to have $\alpha$ of order one, that coal
placement, metallurgical sequence, and the rest are each more likely
than not given the rest of the record, which is a considerably
stronger claim than the objection is usually asked to bear.

\paragraph{Estimating $\alpha$: the sister-case method.} The remaining
empirical input is supplied by an instrument the companion already
develops, and it should be used rather than reconstructed
\cite{Nexus}.

By the Theorem of Aimless Optimation, Optimation is the Cosmic Nexus
conditioned on the attribute $H_N$ that an Aetherian exists in one's
biosphere; and if an attribute $X$ is statistically independent of
$H_N$, the two Nexae assign it the same weight. Independent attributes
are therefore \emph{controls}: their observed frequency estimates the
Cosmic Nexus probability directly, unlifted by conditioning. The
companion's worked case is the Catarrhini node --- Old World monkeys
furnish over a hundred extant species across twenty-five million years
with no attainment of ape-level intelligence, and their fate being
independent of $H_N$, that frequency transfers as the baseline for our
own lineage at the same node.

This is exactly the quantity \eqref{gc:eq:bound} requires. For an
attribute $a_j$ admitting such a control, $\alpha_j$ is read off the
control's frequency, and its contribution to $R_1$ follows.

\paragraph{A corollary on once-only attributes, and the independence
check that governs it.}\label{gc:rem:onceonly}
A corollary of the same theorem bears on attributes observed exactly
once, and it should be read as an inference rather than a
consistency claim. Once a first occurrence has secured what the Golden
Point requires, a \emph{second} occurrence is statistically independent
of $H_N$, it changes nothing about whether an Aetherian exists, and by the theorem it therefore retains its Cosmic Nexus weight and is
not lifted. Optimation supplies what is required and does not
overshoot.

The consequence is a discriminating observation. Write $\lambda$ for
the unconditioned expected count across all opportunities. Under the
hypothesis that $\lambda\ll1$, Optimation must lift the attribute for
$H_N$ to obtain, and lifts it to exactly one; \emph{exactly one
occurrence is what that hypothesis predicts}. Under the hypothesis that
no lift was needed, occurrences are Poisson at rate $\lambda$ and
exactly one is a coincidence of probability $\lambda e^{-\lambda}$. The
likelihood ratio between them is $\left(\lambda
e^{-\lambda}\right)^{-1}$:
\begin{center}\small
\begin{tabular}{@{}rrrrrr@{}}
\toprule
$\lambda$ & $1$ & $3$ & $5$ & $10$ & $20$\\
\midrule
ratio favoring lift & $2.7$ & $6.7$ & $30$ & $2.2\times10^{3}$ &
$2.4\times10^{7}$\\
\bottomrule
\end{tabular}
\end{center}
So a single occurrence tells strongly against the configuration having
been \emph{common} --- $\lambda\ge5$ is disfavored thirtyfold or worse
--- while discriminating only weakly between a genuinely minute
$\lambda$ and one of order unity. Its evidential role is accordingly to
exclude the objection that the prerequisites were unremarkable, not to
supply the magnitude. The magnitude comes from the independent
per-opportunity estimate, and where that gives $\lambda\ll1$ the
attribute contributes approximately $1/\lambda$ to $R_1$.

The corollary is only as good as its individuation, and this is where
the difficulty sits. Rather than proceed directly to an estimate, we
calibrate the instrument on two attributes for which the data are good.
Both return nothing, and the pattern in \emph{how} they return nothing
is what fixes the strategy.

\paragraph{Calibration I: coal.} Britain's coal is the obvious first
candidate, and it fails at the endowment grain. Carboniferous and
younger basins are known on virtually every continent, the
Appalachians, the Ruhr, the Donets, the Kuznetsk, the Nord-Pas de
Calais--Saar--Lorraine field, and others, giving eight or more on any
reasonable individuation, and Britain's is not the most favorable, the
Appalachian seams being unusually large and unusually exposed. With
$\lambda\approx8$ we obtain $\alpha=1-e^{-\lambda}\approx0.9997$ and a
contribution to $R_1$ of about $3\times10^{-4}$: nil.

The verdict is confirmed by the author whose thesis is most at stake.
Allen observes that the Ruhr went unexploited not for want of coal but
because Dutch peat held fuel prices down, removing the return on
improving Rhine transport; and he draws the counterfactual explicitly,
that had German coal been developed three centuries earlier the
industrial revolution might have been a Dutch--German achievement rather
than a British one \cite{Allen09}. He identifies Belgium, around
Li\`ege and Mons, as possessing a comparable ratio of wages to energy
cost. So the claim that the path to the twentieth century ran through
northern Britain was always about the \emph{conjunction} of cheap energy
with dear labor and the institutions to exploit both --- never about
the seams, which were widely available. A reader who recalls that claim
should read it that way, and the estimator's verdict agrees with Allen
rather than contradicting him, reached independently and from geology.


\paragraph{Serial attributes escape the bound.} Proposition
material of this paper's earlier drafts, historical provenance only is a claim about parallel attributes only, and the
restriction is not a technicality --- it identifies the wrong turn that
coal and megafauna represent. An attribute may instead be
\emph{serial}: a conjunction of successive branch points within a single
trajectory, each of which had to go one way for the next to be reached.

\begin{proposition}[Serial attributes; \textbf{[S]}]
\label{gc:prop:cascade}
Let an attribute consist of $m$ successive branch points, the $i$th
resolved favorably with probability $p_i$ conditional on its
predecessors. Then $\alpha=\prod_{i\le m}p_i$ and the contribution to
$R_1$ is $\alpha^{-1}-1$, which grows \emph{exponentially in $m$}. At
$p=0.75$ throughout: $m=8$ contributes about $9$, $m=16$ about $99$,
and $m=25$ about $1.3\times10^{3}$.
\end{proposition}

\noindent The opportunity structure is what differs. A parallel
attribute is bounded by how many sites exist, and there are five
continents. A serial attribute is bounded by how many branch points the
trajectory contains, and a well-articulated historical episode contains
many. Recent history is therefore not weak evidence; it was being
individuated in the form that makes it weak. This is how the candidates should be stated. Not ``Alexander existed'',
which is a parallel attribute of no force, but \emph{that he survived
and succeeded through the entire sequence}, the wound at the
Granicus, Issus, Gaugamela, the near-fatal injury among the Malli, the
absence of a disabling illness before consolidation, and the
Hellenization surviving the immediate fragmentation of the empire, each a branch point with its own $p_i$. Likewise the Roman survival of
the Second Punic War is a sequence, not a fact: the recovery after
Cannae, Hannibal's failure to march, the holding of the Italian allies,
Scipio's African campaign. Likewise the passage of Christianity from a
provincial sect to an imperial religion, and its part in what followed.
Each such episode is one attribute in \eqref{gc:eq:bound}, and each
carries a contribution set by the length of its cascade rather than by
the number of continents.

\begin{remark}[What serial attributes cost]\label{gc:rem:serialcost}
The gain in magnitude is paid for in the independence check, and the
price should be stated plainly, since it is where such an estimate would
fail. For each branch point one must hold that its failure had no adequate
\emph{substitute} --- that had Alexander died at the Granicus, the
Hellenization of the East does not simply proceed under another
commander. History is rich in substitutes, and each $p_i$ that admits
one contributes nothing, since the attribute is then not load-bearing
for $H_N$ at that step. The estimator's discipline accordingly transfers
from grain, where it sat for parallel attributes, to substitutability.

Two consequences follow. Cascade length must be counted in
\emph{substitution-free} branch points, not in dramatic incidents, which
will reduce $m$ considerably below the narrative count. And the
per-step probabilities must be conditional on the predecessors, since
the steps are plainly not independent --- a commander who has survived
four battles is not a random draw at the fifth. Both corrections push
$\alpha$ upward, and a careful estimate should expect the eventual
contribution to fall well short of the naive product.

\begin{remark}[Individuate the branch point by its role, not its
occupant]\label{gc:rem:roleindividuation}
The substitutability audit has a corollary that governs how each branch
point must be \emph{stated}, and getting it wrong is the most likely
route to an inflated estimate. The counterfactual at each step must be taken over the \emph{role}, not
over the occupant. The question at the relevant branch point is not
whether Christianity in particular arose and spread, but whether
\emph{anything} arose that discharged the role it discharged --- in the
window, with the reach, and with the consequences for what followed. So
individuated, the competing cults of the period are not rivals to the
hypothesis but candidate fillers of the same role, and their availability
raises $\alpha$ at that step rather than lowering it: where a role admits
$\kappa$ independent candidates each of probability $q$, the
role-probability is $1-(1-q)^{\kappa}$, not $q$.

Role individuation is therefore the \emph{conservative} choice, and that
is precisely its merit. It concedes the substitution objection in
advance rather than meeting it afterward: a critic who replies that some
other movement would have served is not refuting the estimate but
restating the quantity already being estimated. What survives such an
audit is a residue of steps at which no filler was available with
comparable probability --- and the claim, where a cascade is offered, is
about that residue.

It also relocates where improbability may legitimately be found. It is
not to be sought in the striking identity of the occupant, which is
where narrative attention naturally goes, but in the narrowness of the
window in which \emph{any} occupant could have appeared. The
transitional periods are the interesting ones for that reason, and the
dramatic incidents largely are not.
\end{remark}

\begin{remark}[The limitation is present ignorance, not the structure]
\label{gc:rem:measurement}
One further point about the status of these estimates, and it
cuts in an unusual direction. The obstacle to evaluating $\prod_i p_i$ for any actual cascade is not
that the quantity is ill-defined or that the method is unsound. It is
that counterfactual history is not yet in a state to supply the
$p_i$ with defensible error bars --- the substitutability audits and the
window widths that Remark~\ref{gc:rem:roleindividuation} requires are
matters on which the discipline has views but not measurements. The
uncertainty is \emph{ours}, and contemporary; it is not a property of
the cascades.

Two things follow. The paper should not, and does not, offer a numerical
cascade estimate: what it offers is the structure
\eqref{gc:eq:bound}, the serial--parallel distinction, and the
individuation discipline, with the arithmetic left for inputs that do
not yet exist. And the expectation that such inputs will exist is a
substantive bet rather than a certainty --- one that quantitative
counterfactual history will mature to the point of supplying audited
branch-point probabilities. If it does, the enumeration becomes
routine and the bound follows. If it does not, the weight claim remains
structurally established and empirically unfilled, which is a weaker
position than we would like and a stronger one than assertion
\textbf{[O]}.
\end{remark}

\paragraph{An unnamed significance is a prediction, and it has a
falsification test.}\label{gc:rem:unnamed}
Remark~\ref{gc:rem:measurement} leaves the framework in a position that
looks, at first, like an unfalsifiable auxiliary. We do not know what
the establishment of a universalizing cult in the late-antique
Mediterranean forced through downstream; the analysis nonetheless
concludes that something Golden-relevant was forced through, on the
strength of \ref{gc:prop:flukes} alone. Stated carelessly this is the
canonical bad move --- every fluke acquires an unnamed significance,
nothing counts against the framework, and the Accordance analysis
becomes a device for dignifying whatever happened.

It is not that, and the reason is that the unnamed significance is a
\emph{prediction} with a definite content and a definite refutation. \emph{Content.} For a cascade $C$ satisfying the two admissibility
conditions, genuinely improbable, and load-bearing in the sense of
Remark~\ref{gc:rem:roleindividuation}, the claim is that there exists
a consequence of $C$ on which the probability of the Golden Point
depends, and that it will be identified. The claim is made before the
consequence is known, which is what makes it a prediction rather than a
description.

\emph{Refutation.} Exhibit a route to the Golden Point that does not
require $C$. If the terminus is reachable without the cascade, then
$C$ was not load-bearing for it, and the predicted consequence does not
exist to be found. This is the substitutability audit of
Remark~\ref{gc:rem:serialcost} run at the \emph{terminus} rather than at
an intermediate --- which is the only place it can be run when the
intermediate is unnamed, and which is a harder audit rather than an
easier one, since the chain to be traced is longer.

\emph{And the criterion is diachronic.} A single unnamed significance
is a debt; a growing ledger of them is a symptom. The framework is
progressive if the ratio of identified to unidentified consequences
improves as counterfactual history matures, and degenerating if the
unnamed list only lengthens while the named one does not. We state this
as the test we intend the program to be held to: \emph{if the
historiography improves and the significances remain unnamed, that
counts against the framework and is not to be absorbed by it.} The
passage of time is thereby made a test rather than an excuse, which is
what the position in Remark~\ref{gc:rem:measurement} would otherwise
lack.

One asymmetry should be noted, since it limits what confirmation is
available. Failure to exhibit a $C$-free route is weak evidence at best,
because such routes are counterfactual claims subject to the same
immaturity that motivated the remark. The test is therefore one-sided in
the ordinary way: it can refute a particular significance claim and
cannot establish one, and the framework's standing rests on
whether the identifications accumulate.

\paragraph{The consequence: aggregate, but of serial attributes.} The
comparison requires $R\gg1$, not $R>10^{3}$, and $R_1$ is a \emph{sum}.
Even under the corrections of Remark~\ref{gc:rem:serialcost}, a modest
stockpile suffices: twenty attributes contributing about ten apiece
give $R_1\approx200$, and a single well-attested cascade may supply as
much on its own.

\begin{center}\small
\begin{tabular}{@{}rrr@{}}
\toprule
attributes & each $\approx5$ & each $\approx10$\\
\midrule
$10$ & $50$ & $100$\\
$20$ & $100$ & $200$\\
$40$ & $200$ & $400$\\
\bottomrule
\end{tabular}
\end{center}

\noindent So the program is an enumeration of serial contingencies,
each articulated as a cascade and each audited for substitutability. The
record supplies candidates readily --- the military sequences above; the
demographic rupture of the Black Death and its effect on the very wage
structure Allen's account requires; and the theses of \cite{Diamond97}
beyond megafauna, several of which admit serial rather than parallel
statement. What the two calibrations established is not that recent
history is barren but that it must be interrogated along the trajectory
rather than across the map.

\paragraph{Why the negative calibrations are retained.}\label{gc:rem:whyretain}
Two attributes were examined and both returned nothing. It would be
tidier to omit them, and it would be a mistake. They demonstrate that the instrument has teeth. A method that never
returns nothing cannot discriminate, and a reader entitled to suspect
that the estimator was constructed to deliver a congenial answer is
answered by two cases where it did not. They also fix the grain
discipline \emph{before} any positive estimate is offered, so that the
eventual stockpile is credible rather than convenient. And they preempt
the two objections a reader arrives with, that coal was surely the
scarce thing, and that Diamond's asymmetry surely settles the matter, with the estimator's verdict rather than with assertion.

Most importantly, the pattern in the two failures is itself the
finding. Had one failed and the other succeeded, the first would read as
a red herring. Both failing, and failing for the same reason, is what
establishes material of this paper's earlier drafts, historical provenance only and thereby the
aggregation strategy. The negative results carry the argument. This is where the argument's remaining empirical debt sits, and it is
now specific and bounded. What \eqref{gc:eq:bound} establishes is the
structure: given the $\alpha_j$, the bound follows. What
material of this paper's earlier drafts, historical provenance only establishes is that no single recent
attribute will supply it, so the debt is a stockpile of order twenty
contingencies, each estimated by the sister-case method and each passing
the independence check. The weight claim is accordingly no longer
asserted but \emph{conditionally computed} --- reduced from a
plausibility judgment about an unexamined sum to an enumeration of
bounded terms, each obtained by an instrument calibrated in advance and
each with its required magnitude stated \textbf{[O]}. What remains owed
is the enumeration, and \emph{that} is a research task rather than an
argumentative gap: the terms are individually modest, the sources are
identified, and the standard each must meet is fixed before any is
computed.
\end{remark}

\begin{remark}[The conditioning is not teleological]
\label{gc:rem:teleology}
A reader's first objection will be that conditioning on a future event
to account for present observations is teleology: the Golden Point
reaching backward to arrange its own prerequisites. The framework
answers this at its foundations rather than here, and the answer should
be cited rather than improvised.

The relevant result is the \emph{indiscriminate time principle}, a
corollary of the extrusion principle \cite{Aethus}, whose content is a
pair of mutual non-informations: knowing where the events an attribute's
blanks refer to sit in space and time tells one nothing about when those
blanks resolve for oneself, and knowing when they resolve for oneself
tells one nothing about where they sit. Blankness is a property of the
indexing Aethus, never of a clock reading --- \emph{blanks lack tense},
and an observer's unlearned past is exactly as open as their unformed
future \cite{AethusBrief}.

The objection therefore rests on a premise the framework denies. It
assumes that the Golden Point's futurity makes it a different kind of
object from an unresolved attribute of the past --- that the not-yet-real
is a category. On the Aethic account there is no such category: the
Golden Point is a blank attribute relative to our Aethus, in precisely
the sense that an unread outcome of the deep past is blank, and
conditioning on it is the same operation in both cases. The asymmetry
one feels is real but \emph{entropic} rather than ontological: the
future blocks discriminating information more stringently because it
sits downstream of the entropic arrow, while the past merely happens to
withhold it --- and informational uncertainty is the same thing under
either obstruction \cite{Aethus}.

The mirror-image case matters, because it is a structure the
reader can check without granting anything particular to this work.
Newcomb's problem conditions on a \emph{past} event, the predictor's
sealed placement, which is blank relative to the agent's information
state; and any treatment on which the agent's decision bears on that
placement must decline the assumption at issue here, that a past
occurrence is thereby a settled fact. Past-Alignment conditions on a
\emph{future} event that is blank relative to ours. Same operation,
opposite direction, and by the indiscriminate time principle the
direction is not a difference. A reader who accepts the one and rejects
the other owes an account of which of the two mutual non-informations
they mean to keep.%
\footnote{The companion treatment is \cite{AethusDecision}, which develops
the Newcomb case on exactly this footing: determinacy indexed to an
information state rather than to a position in time, so that an unread past
outcome is a blank entry and blanks lack tense.}

What the Golden Point does here is therefore what any conditioning event
does: it selects, and does not act. The Aetherian Nexus is not
\emph{caused} to contain rocketry by a later passage; the weighting
conditioned on that passage assigns overwhelming weight to histories
containing it, and we read the record we are in. Whether anything
stronger than selection holds, whether the future cone is constrained
rather than merely uninformative, is exactly what
material of this paper's earlier drafts, historical provenance only declines to settle, and the Void hypothesis
is the position on which nothing stronger holds at all.
\end{remark}

\begin{remark}[Why this is not circular, and where it could fail]
\label{gc:rem:circularity}
The obvious objection is circularity: conditioning on the Golden Point
and then finding its prerequisites is no achievement, since the
prerequisites are what the conditioning selects for. The objection would
land against a bare appeal to Past-Alignment, and it is worth being
explicit that it does not land here, because the argument of
\ref{gc:prin:conditioning} is \emph{comparative} rather than
confirmatory. Two candidate conditioning targets are placed side by
side; they issue \emph{different} predictions about the record, one
indifferent to industrial prerequisites, one not, and the record
discriminates between them. A comparison that could have gone the other
way is not a circle.

Three ways it could fail, in decreasing order of seriousness. The weight claim remains the decisive one, though
\S\ref{gc:subsub:weightbound} has reduced it from an assertion to a
single measurement: the bound \eqref{gc:eq:bound} is proved, and one
prerequisite at $\alpha\sim10^{-3}$ suffices, so what would remove the
contradiction is now the specific demonstration that \emph{every}
candidate prerequisite has $\alpha$ of order one. Second, the
argument requires ``observer'' to be robust under exactly these
perturbations; a sufficiently demanding notion of observerhood, one
entailing industrial capacity, say, would collapse the two
conditioning targets together and dissolve the comparison, though at the
cost of making ``observer'' do work it is not usually asked to do.
Third, the comparison establishes that the target is \emph{not}
observer-existence; that it is specifically the Golden Point is
supported by Past-Alignment rather than forced by the contradiction, and
some third target could in principle serve.
\end{remark}


Both arguments are now in hand. The natural-history one carries the
weakness Remark~\ref{gc:rem:circularity} flags, its robustness
premise is contestable, and the human-history one of
\S\ref{gc:sub:humanhistory} does not; a reader who accepts only the
second has accepted the principle, and the first is then the same
inference extended into deeper time.


\section{The Route World: what the strong reading predicts}\label{sec:routeworld}\label{gc:subsub:routeworld}\label{rw:sec}

\begin{quote}\small
material of this paper's earlier drafts, historical provenance only established that the record we hold is the past-cone
restriction of a conditioning on an event in our future cone, and left
open which of three hypotheses governs the forward cone. This section
develops the strongest of them, the \emph{Route World}, in which the
Golden Point stands as a stated attribute of one's Aethus rather than as
a weight over futures, and does so because it is the one hypothesis
with a phenomenology. If it holds, a biosphere approaching its Golden
Point experiences not fewer improbabilities but more, of a specific kind;
that is a claim about the texture of the near future, and unlike almost
everything else in this work, a claim about ours.

Everything here is \textbf{[O]}, with one stated exception: the
conduct argument of material of this paper's earlier drafts, historical provenance only, which is
\textbf{[$\Phi$]} on an adopted valuation and argues for the
hypothesis exactly as little as the rest. The section describes what
the hypothesis predicts, states what would distinguish it from its
rivals, disciplines the historical reading, connects the answer to the
epistemic ceiling, and derives the conduct the open hypothesis leaves
an agent; it argues for the hypothesis nowhere. What it supplies is the framework's most distinctive
prediction, and a reader should not finish this work without knowing the
strong reading has one.
\end{quote}

Void has been given a form it could take
(material of this paper's earlier drafts, historical provenance only), and Weak a statement of what its two
refusals commit it to (material of this paper's earlier drafts, historical provenance only). Strong needs
something different in kind. Void and Weak needed their commitments
drawn out, whereas Strong's difficulty is that it has no empirical handle
at all --- its three metaphysical questions are stated at
material of this paper's earlier drafts, historical provenance only, and none of them is a question any
observation bears on. This section supplies the handle by describing the
world Strong asserts. The description is a consequence of the hypothesis
and not an argument for it.

\begin{definition}[The Route World; \textbf{[O]}]\label{gc:def:routeworld}
The \emph{Route World} is the world of the Strong extrapolation
hypothesis: the world in which the Golden Point stands as a stated
attribute of one's Aethus rather than as a weight over futures, so that
the conditioning is resolved throughout the cone and not merely behind
one. The name is meant in the sense \emph{RNA world} is meant --- a
regime with a characteristic phenomenology, described so that one can
ask whether one is in it.
\end{definition}

\begin{remark}[Whether ``stated'' is coherent for a future-cone
attribute]\label{gc:rem:statedcoherent}
The definition invites an objection that must be met before anything is
built on it. The framework rejects tensed determinacy: an unresolved
future outcome is a blank entry, and blanks lack tense. How can an
observer then hold as \emph{stated} an attribute whose referent lies in
their future cone --- is that not the very error the framework was built
to avoid?

It is not, and the reason is that the objection reads ``stated'' as
``resolved by the clock,'' which is the tensed reading the framework
denies. Statedness is a property of the \emph{index}, of what the
observer's books carry, and not of the referent's date
(Remark~\ref{gc:rem:notense}). An attribute is stated when it is a fixed entry
of the Aethus, whatever its referential event's coordinate; a stated
attribute with a future referent is no stranger than a stated attribute
with an unobserved past referent, and the framework has carried the
latter since its first postulate. What the Strong hypothesis asserts is
that the Golden Point is such an entry for us, fixed in the books,
not merely weighted, and \emph{that} is a substantive claim about
which Aethus we hold, not a category error about tense.

The analytic form of the same point is already in this work.
the technoanatomy companion \cite{TechnoC} distinguishes the interpolation
$\PP^{(t)}$, a family whose conditioning is still being acquired, from
$Q$, a single measure whose conditioning has been absorbed into the base
state --- and identifies the passage to $Q$ with the statement of a
future-referring attribute in present books. The Route World is the
world in which that passage has occurred. Weak is the world of the
family; Strong is the world of the measure; and the difference between
them is exactly the difference between an attribute weighted and an
attribute stated, which the framework can express and does not collapse.
\end{remark}

\subsection{Golden Optimation, and where it differs}\label{gc:sub:goldenopt}

The Route World's
signature is a strengthening of a phenomenon the companion already
names. Optimation reports that a record conditioned on observerhood
displays flukes against the Cosmic Nexus's own rates \cite{Nexus}; but
those flukes lie, at every moment, in the observer's \emph{past} cone,
because the conditioning event, one's own existence, is behind one.
The supply is bounded by what has already happened.

\begin{definition}[Golden Optimation; \textbf{[O]}]\label{gc:def:goldenopt}
\emph{Golden Optimation} is the same phenomenon under a conditioning
event in the future cone. Because the Golden Point lies ahead, the
flukes its occurrence requires need not have occurred yet: there are
moments at which a given required fluke lies in the observer's future
cone, and the observer's moving present passes over it.
\end{definition}

That is the whole differentia and it is sharp. Classical Optimation
never predicts a fluke; it accounts for flukes already in the record.
Golden Optimation predicts them, because the conditioning event has not
yet occurred and the requirements it imposes are still being met. A
Route World is therefore a world in which one keeps encountering
improbabilities \emph{prospectively}, and the rate at which one does so
is the hypothesis's observable content.

\begin{remark}[One phenomenon, and the label is tense-relative]
\label{gc:rem:carryover}
The two Optimations are not two mechanisms with a seam between them.
Take a fluke $F$ required by the Golden Point, occurring at Aethic time
$t_F$, and an observer at $t$. For $t<t_F$ the observer's present has not
reached it and $F$ is Golden Optimation; for $t>t_F$ it stands in the
record and is Optimation in the companion's sense. \emph{Same event, same
conditioning, different label} --- and what changed was the observer's
position, not anything about $F$.

So the label tracks the observer and not the phenomenon, which is
Remark~\ref{gc:rem:notense} applied to Optimation: an observer far enough
downstream reads our Golden Optimation as their past-cone record,
exactly as we read the Cretaceous. The carryover between the two is
therefore smooth by construction rather than by coincidence, and a Route
World contains no moment at which one regime hands over to another.
\end{remark}

\begin{remark}[The symmetry with the companion, and what the corpus is]
\label{gc:rem:rungzero}
One structural fact should be stated in this work rather than left for a
reader to assemble, since it is what tells the reader what the corpus
\emph{is} rather than what either paper is. The companion's inference about the deep past rests on a single licence:
that an unread past outcome is a blank entry, and blanks lack tense, so
that a past fluke is open to inference at all
\cite{Aethus, NexicFoundation}. The companion calls that its
foundational rung, without which no inference about a past attribute can
be drawn. The present work's conditioning on a future event rests on
exactly the same licence: a future outcome is a blank entry, blanks lack
tense, so a future-cone attribute is open to conditioning in the same
sense a past-cone one is (Remark~\ref{gc:rem:notense}, and
Remark~\ref{gc:rem:statedcoherent} for the strong form).

So the two developments are one move run in opposite temporal directions
on one ontology. The companion reads backward from an observer to the
flukes that produced it; this work reads forward from a conditioning
event to the flukes it requires; and the principle that makes either
reading legitimate is the same principle in both. That is not a
coincidence of two papers written by one hand. It is a fact about the
ontology, and it is why the Golden Point is not an extension of Nexic
reasoning but its mirror image --- the same lens turned the other way.
The companion's forward pointer to this work, when it is written, should
be built on that symmetry and on the finding of
\S\ref{gc:sub:humanhistory}: observer-conditioning is exactly right for
the pre-technological record and provably insufficient beyond it, and a
target not indifferent to the technological character of the record
recovers the companion's weighting as its past-cone shadow.
\end{remark}

\begin{definition}[Golden causality]\label{gc:def:goldencausality}
\emph{Golden causality} names the pattern in which an event in the moving
present is accounted for by Golden Optimation, by the Aethus's
conditioning on the Golden Point, rather than by the Cosmic Nexus's
own rates.
\end{definition}

The name is admissible for a reason internal to the framework rather than
by indulgence. Causal talk here is not talk of a primitive relation of
production: this work's causal vocabulary labels patterns of conditioning
and dependence in the Aethic structure and takes no position on any
further relation underneath, because the framework posits none. Where
causation is not fundamental, a \emph{kind} of causation is a kind of
conditioning pattern, and Golden causality is one --- named, and carrying
no more than Theorem~\ref{imp:targeted} licenses. What the name must not
smuggle is the register Remark~\ref{gc:rem:teleology} forbids: no event
pursues the Golden Point, and Golden causality describes what accounts
for something, never what something is for.

\begin{remark}[Invested Golden causality, as phenomenology]
\label{gc:rem:invested}
One element of the phenomenology deserves its name here, because it is
the form in which Golden causality is most naturally read off a record
from inside: the register of \emph{investment}, ``look how much the
conditioning invested in this; it must have been necessary.'' The
register is earned rather than decorative:
the retrocausal-accordance material set aside for a further paper is its exact content, an observed
route's improbability lower-bounding what it accomplished at a stated
exceedance confidence, and the retrocausal-accordance material set aside for a further paper is why the
economic vocabulary of the decision-theoretic material of the broader program, set aside here was built to carry it, Aurutance's two terms are the cost and the necessity, and the
investment reading is the decision rule run backward. What the register
must never smuggle is agency: nothing takes the cheaper route or
declines it (the retrocausal-accordance material set aside for a further paper), and the expenditure is the
conditioning's concentration of weight, never anyone's spending. A
Route World is, in this register, a world whose record reads as heavily
and legibly invested --- and \S\ref{gc:sub:historicalcases} is where
the temptation to read particular centuries that way is disciplined.
\end{remark}

\begin{remark}[The future-conditioned tradition, placed]
\label{rw:rem:futurecond}
The Strong reading has a published ancestry, and the placement
discipline owes it a paragraph: Barrow and Tipler's Final Anthropic
Principle \cite{BarrowTipler86}, Tipler's Omega Point program
\cite{Tipler94}, Leslie's axiarchism \cite{Leslie79}, and,
ancestrally, Teilhard \cite{Teilhard55} --- the branch of the
anthropic tradition that conditions on the \emph{future}, which is
the branch this work lives on (the companion engages the tradition's
observer-conditioning trunk; the future-conditioning branch is placed
here). The deltas are real and they all cut one way. The conditioning
of this work is a \emph{measure operation}, not a dynamics: no agent
at any grade, no attractor, no physical boundary condition --- where
the Omega program is a physical eschatology with cosmological
commitments, several since falsified. This work carries published
falsifiers and an epistemic ceiling that \emph{forbids} the
certified-salvation register the Omega program embraced
(Theorem~\ref{imp:ceiling}); and it is severable, under Void the
mathematics and the critique stand with nothing conditioned on, where the tradition's packages fall whole. The delta that flatters
neither side: this framework purchases strictly less, no
resurrection mechanics, no convergence guarantee, no collective
eschatology, and prices what it does purchase.
\end{remark}

\subsection{The empirical signature, and its limits}\label{gc:sub:credencemethod}

Two things about the credence being computed here should be fixed before
the computation, since they determine what the computation is a
computation \emph{of}. \emph{The credence is methodological, and it is a superposition.} A
reasoner in a Route World holds the Golden Point as a stated attribute
--- that is what the hypothesis says of their Aethus. But the reasoner's
\emph{credence} that they are in a Route World is computed from less
than their Aethus contains, since it is computed from the evidence
accessible to them and not from the stated attribute itself. So the
credence is an Aethic disagreeing superposition: the Strong hypothesis's
determinate conditioning weighed against a Weak-type alternative in which
the Golden Point is only weighted. The distinction is between
methodological credence --- what the evidence supports --- and
ontological credence --- what one's Aethus in fact carries --- and
everything below is about the first, and never asserts the second.

\emph{The procedure does not force the conclusion.} What it establishes,
if the flukes keep arriving, is that each next event acts \emph{as if}
the Golden Point were being conditioned on, to whatever precision the
accumulated evidence supports. That is the whole of the method's
ambition, and it is the correct one: a reasoner genuinely in a Route
World tends toward the conditioning by this route with arbitrarily high
credence, and never reaches it, and the ceiling of
Theorem~\ref{imp:ceiling} is why the second clause is not a defect.

The Route World is testable in
principle, and the form the test takes is the subsubsection's main
result.

\begin{proposition}[Credence accumulation; \textbf{[P]} given the
premises]\label{gc:prop:credence}
Let $E$ be the conjunction of fluke attributes observed over an interval
of Aethic time, with $q=\PP(E)$ their weight under the unconditioned
law and $Q(E)\approx1$ under the survivorship-conditioned law because
the target requires them. The comparison is between measures, not
between $G$ and $\neg G$, since in the critical regime conditioning on
$G$ as an event does not exist:
\begin{equation}\label{gc:eq:bayesq}
\mathrm{BF}_{Q:\PP}(E)=\frac{Q(E)}{\PP(E)}\;\approx\;\frac1q,
\qquad\text{so}\qquad
\log\text{-evidence}\;\approx\;-\log q .
\end{equation}
For finite information $\Fc_t$ the factor is the density itself,
$\mathrm{BF}_t=dQ|_{\Fc_t}/d\PP|_{\Fc_t}$, written generically as
$Z_t$; in the hard-wall model of Section~\ref{sec:imports} it is the Doob
density the Doob density of \cite{SurvA}, $Q^{\mathrm{hard}}$, proved, while the
physically intended soft law $Q^{\mathrm{soft}}$ has density
$e^{-\int_0^tq}\,g(X_t)/g(X_0)$ conditional on Ansatz~B of \cite{SurvA},
and every result below that consumes the density says which of the two
it uses: the
realized log-evidence is $\ell_t=\log\hm(X_t)-\log\hm(X_0)$, and the
historical Bayes-factor argument and the stochastic $\hm$-transform are
one object. The realized $\ell_t$ is the logarithm of a
$\PP$-martingale and fluctuates along a path; what rises is its
$Q$-expectation, the information of the information coordinate of \cite{TargetedB1}.
\end{proposition}

Three things follow, and the second is the one that keeps the procedure
honest. \emph{The evidence is the same object the decision-theoretic material of the broader program, set aside here computes.} $\hm$ is
what the Golden Score is a ratio of (the Golden Score of the broader program, set aside here),
so the Route World's credence and the decision rule consume one
quantity, and both wait on the same missing map
(the identification map, an open problem of the broader program). \emph{And that map is what stops the procedure being unfalsifiable.}
Absent it, ``what counts as a fluke'' is chosen after the fact, and a
sufficiently motivated observer finds flukes anywhere --- the
look-elsewhere problem, which would make \eqref{gc:eq:bayesq} an
instrument for confirming whatever it was pointed at. With it, a fluke
is whatever raises $\hm$, which the model fixes and the observer does
not. The methodology's rigor is therefore hostage to
the identification map, an open problem of the broader program's open problem, and should be described that way
rather than advertised as available.

\emph{The aggregate does not determine its own decomposition, and
identifiability is a matter of degree.} The likelihood ratio is a
function of the \emph{state}, $\hm(X_t)/\hm(X_0)$, and not a ledger of
which happenings supplied it: an accumulated $20$ nats, a Bayes factor
near $5\times10^{8}$, is equally the signature of one unmistakable
event, of two hundred each amounting to a ten-percent tilt, or of two
thousand at one percent apiece, and the aggregate does not distinguish
the cases.

It does not follow that no individual event is identifiable, and the
qualification matters. A happening of conspicuously low weight under the
Cosmic Nexus is exactly the case
Theorem~\ref{imp:targeted} speaks to: its update factor on
futures is correspondingly large, so a future accounting for it is
favored in the stated ratio, and the principle is thereby the
\emph{detection instrument} for such events rather than merely a
constraint on them. Large single contributions can be found, and an
observer who knows to look will find some.

What resists is the rest, and the shape of the resistance is steep.
Writing $\theta$ for the weight at which a single event becomes
conspicuous and $\mu$ for the typical contribution, the share of total
evidence residing in conspicuous events falls like $e^{-\theta/\mu}$: at
$\theta$ a twenty-to-one improbability and a total of $20$ nats, three
large events put ninety-two percent of the evidence in plain sight,
twenty events put twenty percent, and two hundred put
none.\footnote{Exponential contribution sizes at fixed mean, for which
the conspicuous share is $(1+\theta/\mu)e^{-\theta/\mu}$ exactly.} And
under the log-spread laws this work uses elsewhere,
(the technoanatomy companion \cite{TechnoC}), the two quantities that matter come
apart: with contributions lognormal at spread $3$ over two hundred
events, roughly \emph{one} carries sixty-three percent of the evidence.
A handful is findable, those few are decisive, and the proportion of
contributing events they represent is under one percent.

So the accurate statement is that a Route World should be expected to leave
a few conspicuous traces and a great many invisible ones, the conspicuous
few carrying a large share of the evidence and a negligible share of the
count. Evidential dominance and numerical proportion are different
quantities, exactly as the technoanatomy companion \cite{TechnoC} finds them for
orderings, and an observer who located every conspicuous trace would
still not have located most of what happened. That is
the broader Golden program, set aside here's class-versus-instance discipline arriving from the
evidential side.

Three tasks are worth separating here, because the paragraph above runs
two of them together and the third is the one with a future.
\emph{Detecting} the aggregate needs no localization whatever:
\eqref{gc:eq:bayesq} reads the accumulated ratio off the state, and two
hundred invisible ten-percent tilts deliver the same decisive
$5\times10^{8}$ as one conspicuous event would. \emph{Localizing the
tail} is what Theorem~\ref{imp:targeted} does. \emph{Localizing
the bulk} has no instrument at all, and the reason is specific: the
accordance principle is a \emph{marginal} test, it interrogates one
attribute's weight under the Cosmic Nexus and infers from that weight
alone, so an event contributing at a magnitude no marginal test
resolves is invisible \emph{to it}, which is a fact about the instrument
and not about the event.

That distinction leaves room, and the room should be marked rather than
conceded. The accordance principle is the elementary instrument here: it
requires only a weight and a comparison, it is available immediately, and
it is correspondingly coarse. Nothing establishes that it is the best
available in principle. What a finer instrument would have to do is
compute an event's contribution to $\hm$ \emph{without routing through
that event's marginal improbability} --- and an instrument of that
description is not hypothetical, it is
the identification map, an open problem of the broader program's missing identification map, which is exactly a
rule assigning states to events by their own structure rather than by
their rarity. The Route World therefore acquires the same open problem
the decision-theoretic material of the broader program, set aside here has, from the other direction, and a second reason to want
it solved: the map would convert localization from a tail-only
capability into a general one. Until then, the coarse instrument is what
there is, and the estimate that a handful of traces is findable while the
bulk is not should be read as an estimate about present instruments
rather than a claim about what a Route World conceals \textbf{[O]}.

\begin{remark}[The attribution is lost in the representation, not in the
world]\label{gc:rem:flattening}
It is worth being exact about where the localization problem comes from,
because the answer says what a better instrument would have to do rather
than merely that one is wanted. The Aethic object is a structure: attributes standing in dependency
relations, with weights over the fan. The model of Section~\ref{sec:imports}
projects that structure onto a scalar, a hazard, and downstream a
single harmonic value, and the projection is many-to-one by a wide
margin. \emph{The attribution information is discarded at the
projection.} Nothing about the world hides which events contributed;
the representation this work computes in does not carry the answer,
because a number cannot. That is why the identification map is hard: it
is an inverse to a projection that discarded what the inverse needs, and
inverses of that kind are recovered by supplying structure rather than by
supplying data.

It also says what the accordance principle is finding. A marginal test on
weights locates the attributes of high dependency degree, the ones on
which much else turns, and which are correspondingly improbable in
isolation, so what the principle recovers is a \emph{skeleton}: the
trunk and the hubs of the dependency graph, and not its edges. Working in
the unflattened structure would be working with the graph rather than its
scalar summary, and there the bulk contributions are not small numbers
lost in an aggregate but edges with endpoints. Whether the framework's
dependency structure supports inference at that grain, and what
volume of recorded attribute-data would make the closure computable, is
open, and is a different open problem from the identification map rather
than the same one \textbf{[O]}.
\end{remark}

\begin{remark}[The collapse condition is also a legibility condition]
\label{gc:rem:legibility}
One consequence runs the other way and bears on the companion rather than
on this work. The companion's collapse condition,
$c_{\mathrm{eff}}>\log_{10}(n(b-1))$, is stated as a fact about
\emph{weight}: above it, the realized ordering carries a share of order
unity. It is equally a fact about \emph{inference}, and the second
reading has not been drawn.

Take the effective number of live alternatives to be the exponential of
the entropy over the tier structure. At twenty interrogable slots of ten
values the nominal space is $10^{20}$; at $c_{\mathrm{eff}}$ just above
the threshold it is a few hundred; at $c_{\mathrm{eff}}=3$ it is four;
and by $c_{\mathrm{eff}}=5$ it is one to within a percent. The collapse
is sharp, and what collapses is the space a reconstruction must search.
So the same contrast that makes the record improbable makes it
\emph{legible}: reconstruction of a high-contrast natural history
requires data scaling with the surviving handful rather than with the
nominal branching, which is the precise content of the intuition that
knowledge here should grow quickly once methods are good enough to
exploit the concentration. The threshold such methods must reach is the
collapse condition itself.

Two cautions. The concentration is \emph{conditional}, it is what the
conditioning purchases, and without it the space is the nominal one, so the legibility cannot be cited as unconditional evidence for the
conditioning without circularity. And concentration is not determinism:
one ordering carrying a share of order unity leaves the others weighted
rather than excluded, and Definition~\ref{gc:def:epsdep}'s grading is exactly the
difference. What is available honestly is a prediction: if the
conditioning holds, reconstructions of deep history should converge on
less evidence than a flat-prior analysis of the same branching would
require, and that is checkable without any fluke being counted
\textbf{[O]}.
\end{remark}

One further consequence cuts against the methodology's own convenience.
How large the required flukes are depends on how much work the
conditioning has to do: a position from which the Golden Point is already comparatively
likely requires only small corrections, while deep time, from which it
was not, required the enormous ones the record displays. A biosphere
improving its position should therefore expect its Golden-Optimation
events to grow \emph{less} conspicuous, so the evidence available to
\eqref{gc:eq:bayesq} thins exactly as the situation it reports on
improves, and the method is at its most informative where the news is
worst \textbf{[O]}.

\emph{The ceiling permits the ambition and forbids the conclusion.} By
Theorem~\ref{imp:ceiling} the two measures are mutually absolutely
continuous on every finite window and singular only on the
infinite-horizon $\sigma$-algebra, so the likelihood ratio is finite at
every finite time: credence may rise without bound and \emph{never
certifies}. That is not a limitation discovered here but the theorem's
own shape, and it is the exact sense in which the procedure can approach
the conclusion asymptotically while being barred from reaching it.


\subsection{What would distinguish the three worlds}
\label{gc:sub:discriminate}

A hypothesis with a phenomenology owes a statement of what would tell it
from its rivals, and this is where the Route World's status is decided.
The answer has two halves and they point in opposite directions.

\paragraph{What differs.} The three worlds make different predictions
about the \emph{rate} at which an observer encounters improbabilities
against the Cosmic Nexus's own rates.

\emph{Void.} Nothing constrains the forward cone, so improbabilities
arrive at the Cosmic rate: after one's present, flukes are as rare as the
unconditioned weighting says, and Optimation is a fact about the record
only, and the realized log-likelihood favors $\PP$ in expectation:
$\gamma_\PP<0$. \emph{Weak.} The forward cone is weighted by Golden-Point probability, so
improbabilities arrive at a rate above Cosmic and below what full
conditioning gives --- and, since Weak's anchor does real work, at a rate
depending on one's position. No rate prediction exists for Weak until
its law is specified: a partial-conditioning family $R_\theta$, a
tempered Doob transform, or a finite-horizon conditioning regime, with
$\gamma(\theta)=\lim_t\E_{R_\theta}[\ell_t]/t$ computed and the
ordering $\gamma(0)<\gamma(\theta)<\gamma(1)$ proved, which is open
\textbf{[O]}.

\emph{Strong.} The conditioning is resolved throughout the cone, so
required improbabilities arrive at the full conditioned rate, and by
Proposition~\ref{gc:prop:credence} the realized log-likelihood favors
$Q$ in expectation at the $Q$-expected rate, $\gamma_Q=\lim I_t/t>0$, while
$\ell_t$ fluctuates about its mean. The discriminating quantity is therefore the sign of the expected
log-likelihood growth under the actual law: Void and Strong predict
opposite signs for the $Q{:}\PP$ log-evidence under their own
generating laws, and Weak awaits its law. Two things bear on reading it.
Remark~\ref{gc:rem:carryover} says the label attaching to any given fluke
is tense-relative, so an observer cannot distinguish the worlds by
inspecting individual events; only the \emph{rate} discriminates. And by
the identifiability analysis of \S\ref{gc:sub:credencemethod}, the rate
must be read from the aggregate $\Lc_t$ rather than from a ledger of
attributed events, since the aggregate does not decompose. What one
compares is a log-likelihood ratio against a Cosmic-rate baseline: Void
predicts that it falls in expectation, at the rate $-D(\PP\|Q)$, and
Strong that it grows.

\paragraph{What does not.} By Theorem~\ref{imp:ceiling}, $Q$ and $\PP$
are mutually absolutely continuous on every finite window and singular
only on the infinite-horizon $\sigma$-algebra. So the discriminating
quantity is finite at every finite time. Evidence for Strong can
accumulate without bound and never certify; evidence for Void can only
ever be the absence of accumulation, which is compatible with Weak at low
flux and with Strong at early time. \emph{No observation of bounded
duration decides among the three.}

That is not a limitation the Route World discovers about itself but the
ceiling's own shape, and it should be read as the point rather than as a
disappointment. What the framework predicts is a \emph{trend}, a
log-likelihood ratio that rises in expectation under Strong and falls in
expectation under Void, Weak having no rate prediction until an explicit
intermediate law is supplied, and a trend is exactly the kind of
prediction that accumulates credence without converting to certainty. A
reader who wants a decisive test is asking for something the framework
proves it cannot supply; a reader who wants a discriminating quantity has
one.

\paragraph{The scoring.} The Route World therefore has a phenomenology, a
discriminating quantity, and no decisive test --- and the last is a
theorem. That is a weaker empirical status than a falsifiable prediction
and a stronger one than a metaphysical hypothesis with no observational
content, which is where Strong stood before this section. What the
section has done is move Strong from the second category into the first,
and no further.

\begin{remark}[Prepositioned prerequisites: a second discriminating
quantity, physical rather than historical; \textbf{[S]} with an open specification problem]
\label{rw:rem:preposition}
The flux is a historical quantity, read from the record of events
that have already occurred. The author proposes a second
discriminating quantity of a different kind, read from physical
structure that has not yet been used. Suppose the route to the
Golden Point requires resources this planetary system either has or
lacks, of the sort a departure from the system would consume, and
suppose those resources are not required for anything that has
happened so far. Under the strong reading the conditioning is
resolved throughout the cone and the route is realized, so whatever
the route needs must be in place; the requirement is therefore
already satisfied in structure standing today and available for
inspection, well before anyone consumes it. The author's
illustrative case is water ice in outer-system locations that later
turn out to matter, and the example is illustrative only, carrying
no weight of its own.

The reason this quantity is worth stating separately is that it
escapes the blur that limits the first. Remark~\ref{gc:rem:carryover}
observes that the label attached to any individual fluke is
tense-relative, so an observer cannot sort worlds by inspecting
events, only rates. An \emph{unconsumed} prerequisite is the one
class of feature to which that objection does not apply, because
nothing in the past has used it: it cannot have been
observer-necessary, having so far been necessary for nothing at all.
This makes prepositioning a genuine discriminator between the
readings rather than a further instance of the same evidence, since
observer conditioning constrains only what the past required, and
says nothing whatever about what the future will need. Two things
must be controlled before the discrimination is real. The resource
must not be a byproduct of the same processes that produced
habitability, which is the common-cause confound, and the base rate
must be computed conditional on habitability so that observer
selection is not doing the work unnoticed. With those controls the
evidential structure is exactly that of
Proposition~\ref{gc:prop:flukes}: the likelihood ratio is one over
the probability that an arbitrary habitable system satisfies the
stated requirements, and the force of the test is the smallness of
that probability, which comparative planetology may eventually
estimate as the census of systems grows. The falsifier follows
directly. If our system's provisioning against requirements
specified in advance turns out median or worse among habitable
systems, that counts against the strong reading rather than for it.

The specification problem is the honest catch, and it is the
author's own. The test has force only when the requirement is fixed
before the provisioning is inspected, or the retrospective-target
null of Remark~\ref{gc:rem:sharpshooter} consumes it entirely, and
fixing the requirement demands joint knowledge of what the route
will need and of what supplies it, which a species at our epistemic
stage does not have. The conservative form of the claim is therefore
addressed to descendants rather than to us: as the epistemology of
Golden-Point necessity matures, prerequisites should be found to
have been visible in the physical record long before their necessity
was understood, so that the sentence available to our descendants is
that a feature already lay open to our inspection and turned out to
be required. That is a prediction about a future epistemic state,
scoreable by those who reach it and not by us, and it is entered at
the same grade and with the same honesty as the reception-scale
predictions of the decision-theoretic material of the broader program, set aside here. The ceiling governs throughout, since
a satisfied requirement raises the ratio and certifies nothing.
\end{remark}

\subsection{Why the historical cases do not discriminate}
\label{gc:sub:historicalcases}

A reader who accepts the previous subsection's discriminating quantity
will nonetheless want to reach for something nearer to hand, and the
material is genuinely there: read a consequential century closely and
what had to go right is startling. This subsection says what such cases
do and do not support, because the temptation to over-read them is
strong and the framework's own guards are what keep the reading honest.

\paragraph{The material.} If one is in a Route World then one has been in
it throughout, and any century should show it. The scientific
enlightenment is the case most often reached for, and it repays the
reach: the transmission of algebra; the printing press and the literate
market it made \cite{Eisenstein79}, and --- less often noticed and more interesting --- a
theological temperament under which fitting exact laws to the world was
a natural thing to attempt rather than an eccentric one --- a standing
thesis of the historiography \cite{Merton38}. Counterfactuals
of the usual shape suggest themselves. Had Rome persisted in a form that
left no room for the religious configuration that followed, it is not
obvious that the temperament distinguishing Renaissance mathematics from
Archimedean mathematics would have arisen; and the material prerequisites
have their own contingencies, paper among them. Each century then sets
conditions for the next, which is the schedule structure of
the schedule of \cite{TechnoC} in a documentary register: gates in order, with
the ordering doing work.

\paragraph{What they support.} The cases are the human-history leg's own
material (\S\ref{gc:sub:humanhistory}), and what they support is what
that leg argues: that the conditioning is sensitive to the technological
character of the record, which observer-conditioning is not. That is a
claim about \emph{which target}, and these cases bear on it.

\paragraph{What they do not.} They are past-cone cases, and the
distinction that matters is exactly the one
Definition~\ref{gc:def:goldenopt} draws. A fluke already in the record is
\emph{consistent with} Golden Optimation and is equally consistent with
ordinary Optimation --- the conditioning event behind one, the supply
bounded by what has happened. By Remark~\ref{gc:rem:carryover} the label
attaching to any past fluke is tense-relative, and a fluke that lies
behind an observer carries no information about whether the conditioning
extends ahead of them. So historical cases are \emph{Past-Alignment
consistent}, which is the milder claim, and they do not discriminate
Strong from Weak or from Void. Only the flux does, and the flux is a rate
in the observer's forward cone.

The stronger phrasing, \emph{documentable cases of Golden causality}, should accordingly be resisted. By
Definition~\ref{gc:def:goldencausality} Golden causality names events
accounted for by the Golden conditioning \emph{rather than} by Cosmic
rates, and establishing that of a particular event requires establishing
the conditioning, which is what is at issue. What is documentable is
consistency, and consistency with a hypothesis is not evidence for it
over a rival the same data also fit.

\begin{remark}[Three guards, and the one that bites hardest]
\label{gc:rem:historyguards}
The historical material is where this work's own cautions apply with
most force, and they should be applied before the material is used
rather than after it is challenged.

\emph{The scan is enormous.} material of this paper's earlier drafts, historical provenance only's multiplicity
guard bounds expected $\kappa$-fold violations among $n$ interrogated
pairs by $n/(1+\kappa)$, with the standing condition that $n$ counts
attributes \emph{scanned} and not merely those interrogated. For a
century read by historians, the scan is every attribute anyone has ever
found notable, and an honest $n$ is enormous. The noise floor is
correspondingly high, and a historical trace must stand clear of it by a
wide margin before it is worth anything. This is the guard that bites
hardest and it is the one most easily skipped.

\emph{Substitution is the question, not occurrence.}
\S\ref{gc:sub:humanhistory} audits by \emph{role} and not by occupant:
the question is never whether Christianity in particular arose, but
whether anything arose that discharged the role. Applied to the
temperament case this cuts both ways, and interestingly. Material
prerequisites look substitutable --- many routes lead to mass literacy,
and the printing press is one occupant of that role. A specific epistemic
disposition looks less so, since it is not obvious what else would have
discharged it; and if that holds, temperament is a \emph{serial}
attribute in the sense of Remark~\ref{gc:rem:serialcost} where the technology
is parallel, and contributes correspondingly more to $R_1$. That is a
conjecture about substitutability and not a finding, and it is exactly
the kind of claim the serial-attribute programme exists to settle
\textbf{[O]}.

\emph{And narrative bias is prior.} Remark~\ref{gc:rem:narrativebias}
states the objection and the asymmetry that answers it, and nothing in
this subsection escapes it: dramatic near-misses make better stories and
are more often recorded, so a century dense with them may be reporting a
fact about historiography. The test remains the connectivity asymmetry
that remark names, and no accumulation of vivid cases substitutes for it.
\end{remark}


\section{What is not claimed}\label{sec:fences}

\subsection{Fences}\label{gc:sub:routefences}

\begin{remark}[Determinate is not accessible, and why the distinction
cannot be relaxed here]\label{gc:rem:determinacygap}
An objection is available against the fences below, and it is good enough
that stating it is better than hoping no reader finds it. It runs as a
model comparison. \emph{The objection.} The fences say that whether one's own contribution
is among the conditioning's requirements is a thing no observer can
establish. But the Route World is precisely the hypothesis that the
conditioning is \emph{resolved} --- the Golden Point stated rather than
weighted, throughout the cone. If it is resolved, the requirements are
settled; if they are settled, there is a fact about whether any given
contribution is among them. So either that fact is available, in which
case the fences fail and the complacency inference goes through --- or it
is not available even in the resolved case, and one wonders what
``resolved'' was supposed to mean. Push the first horn and the Route
World undermines the conduct it recommends; push the second and it looks
like a hypothesis with no content. Either way, so the objection
concludes, the Route World is incoherent and its incoherence is a proof
that we are not in one.

\emph{The answer, and it is not a concession.} The objection equivocates
on ``available.'' Statedness is a property of an \emph{Aethus}, of
what the world's books carry, and the framework is explicit that an
Aethus is deductively closed and indexed without regard to whether its
bearer can compute over it \cite{Aethus}. So in a Route World the
requirements are determinate \emph{in the books} and remain inaccessible
\emph{to the reasoner}, and no contradiction arises, because those are
different predicates of different objects. What
Theorem~\ref{imp:ceiling} forbids is not there being a fact but
\emph{certification of it from observation}: $Q$ and $\PP$ are mutually
absolutely continuous on every finite window, so no finite record
distinguishes them. That is a limit on inference and not on determinacy.

\emph{Two corrections the objection earns.} First, the fences should not
have said \emph{in principle}, and above they no longer do: the measures
are singular on the infinite-horizon $\sigma$-algebra, so an observer of
unbounded duration would distinguish them, and what the ceiling supplies
is unavailability to any \emph{finite} observer. Second, the resolved
case is not contentless. What ``resolved'' means is
the technoanatomy companion \cite{TechnoC}'s distinction, one is under the
measure $Q$ rather than the family $\PP^{(t)}$, and that has
consequences the section spends its length on, including the evidence
flux that orders the three worlds. A hypothesis can be substantive and
its truth still uncertifiable; this work's own epistemic ceiling is the
standing example, and it applies to this work's conclusions as firmly as
to its targets.

\emph{And the structure is one the section already carries.}
Remark~\ref{gc:rem:flattening} makes the same distinction on the
evidential side: the attribution information is discarded \emph{at the
projection}, not in the world, so which events supplied the accumulated
evidence is determinate and not recoverable. Determinate-and-inaccessible
is the framework's normal condition rather than a special pleading
introduced to save the fences. \emph{And a route that closes the leverage without the diagnosis.}
Even were the objector granted the availability premise arguendo, no
model comparison against the Route World follows, for a structural reason
independent of the equivocation above: the ceiling is a property of the
model under every hypothesis of \S\ref{gc:sub:discriminate}, Void,
Weak, and Strong alike, so a premise that falsifies all three worlds
at once selects among none of them. What an exhibited certification
would refute is the mutual-absolute-continuity structure itself, at the
site where its empirical exposure lives (the survivorship companion \cite{SurvA}), and
the Route World would stand or fall with its rivals rather than before
them.
\end{remark}

\begin{remark}[The robustness argument, and why it is not relied on]
\label{gc:rem:routerobust}
A stronger conclusion suggests itself and should be marked as the
weakest thing here. If each conservative alternative, conditioning on
some nearer, more peripheral future event rather than on the Golden
Point, is one whose own passage will dissolve whatever objectivity it
seemed to have, then only a target at the asymptote remains stable under
the motion the argument itself describes; and one may be tempted to
conclude that nothing \emph{but} the Golden Point can be coherently
conditioned on.

Two reasons not to lean on it. It requires that the peripheral
alternatives be exhaustively surveyable, which is the same
exhaustiveness the companion's own arguments purchase carefully and this
one assumes. And it is the argument's only step that is not a
likelihood computation, so a reader who accepts \eqref{gc:eq:bayesq} and
Theorem~\ref{imp:targeted} and rejects this paragraph loses nothing else.
The accumulation is what gives the Route World teeth; the robustness
claim is a conjecture stated beside it \textbf{[O]}.
\end{remark}

\begin{remark}[What the Route World does not license]
\label{gc:rem:routefence}
One reading has to be closed explicitly, because it is the natural one
and it is dangerous. If required flukes occur, then, the reading goes, an individual
contemplating some improbable feat should ask not whether they will
succeed but whether the Golden Point requires their succeeding, and
proceed accordingly. Nothing in this work supports that, and two of its
results forbid it. By the broader Golden program, set aside here an individual is not a
lever on their biosphere's hazard: what they hold is a position, not a
purchase, and Golden Optimation is a statement about which histories one
finds oneself in rather than about what any agent can bring about. And
by Theorem~\ref{imp:ceiling} nobody can determine whether their own
feat is among the requirements, since no finite observation places a
system on the interpolation at all. The premise of the reading is therefore
unavailable to any finite observer, which is not the same as there being
no fact of the matter (Remark~\ref{gc:rem:determinacygap}).

What survives is weaker and is the defensible form: a Route World is one in
which improbabilities are more common than the Cosmic Nexus's rates
predict. It is not one in which any particular improbability is owed to
any particular person, and an agent who acts on the stronger reading is
not applying this framework but discarding
Theorem~\ref{imp:ceiling}. the broader Golden program, set aside here fences the neighboring
case in the same terms and for the same reason.
\end{remark}

\paragraph{The phenomenology extended: from events to constitution.}\label{gc:rem:constitution}
The Route World as developed in this section is a claim about
\emph{events}: which occurrences carry weight, how the record reads
under the conditioning, what the forward flux does. Its Strong reading
supports one further extension, developed where the categorical
machinery lives rather than here: from the structure of events to the
constitution of things --- the universe's contents as
weight-carrying excursions on the Aethus that states the conditioning
attribute (material of this paper's earlier drafts, historical provenance only). Nothing about the Route
World's own claims strengthens or weakens by the extension, and its
fences govern there unweakened; the pointer is recorded here because a
reader of this section should know the hypothesis's full reach before
weighing it.

\begin{remark}[A consequence for periodization]\label{gc:rem:periodization}
One application is worth recording because it is where the Route World
meets Principle~\ref{gc:prin:continuity} and returns something concrete. Golden Optimation running over the ancestral record and over the present
is one process under one conditioning, so there is no categorical
transition between the environments of natural history and those of the
present --- which is continuity's own claim, arrived at from the
conditioning side. The present is nonetheless distinguished in
\emph{stakes}, by the broader Golden program, set aside here, and the two are
compatible for the reason the broader Golden program, set aside here gives: a
monotone sequence distinguishes every later point over every earlier one
without distinguishing any point categorically. The right image is a
field potential that is high here because the Golden Point is near ahead
--- and the misreading to guard against is that a high potential marks a
step. It does not: a potential rises continuously toward its source, and
what is high at the present is high because of proximity to a transition
that lies ahead, not because a transition has been crossed.

The consequence for periodization is that a boundary drawn in the past, an Anthropocene, on any of its proposed placements, marks a
change in stakes rather than in kind, and continuity is what denies it
the second reading. The argument should be run from the principle and
not from the failure to locate a line, since inability to place a
boundary is a sorites and not a proof. What the Route World adds is a
positive alternative: the transition the periodization is reaching for
lies \emph{ahead}, at the Golden Point, which is the one place the
framework does mark a change in kind --- and which
Theorem~\ref{imp:ceiling} forbids anyone from dating.
\end{remark}



\subsection{What this section does not claim}\label{gr:sub:fence}

The fence is stated hard, on the companion's discipline that an
apparatus which prices what it claims must also refuse what it does not
\cite{Ben26Optimal}. It is not claimed that any observed system is Golden, ideally or
practically; by Theorem~\ref{imp:ceiling} the ideal version is
unknowable to any finite observer (the technoanatomy companion \cite{TechnoC}) and
the practical version is out of reach for want of an ensemble. It is not claimed that the companion's cosmic
weighting and the path measure $\PP$ of
Section~\ref{sec:imports} are the same object; that
identification is flagged as unperformed in
the technoanatomy companion \cite{TechnoC} and nothing here depends on it. It is not
claimed that any individual biosphere, ours included, will behave as
$the technoanatomy companion \cite{TechnoC}$ describes; those are statements about a
conditional distribution. No normative conclusion is claimed
($\Phi2$). And the companion's natural-historical program, its
conjectures, its bounds, its paleontological predictions, is not
claimed, imported, or re-argued; $\Phi5$ enters as a stated result and
nowhere else.

One thing remains on the far side of the fence, because a fence is
legible only if what it excludes can be seen. A program built on this
section would inherit three commitments, stated here as refusable rather
than as claims. It would be committed to the boundary-cost scaling of
$\Phi3$ showing up in institutional and ecological data, not merely in
argument. It would be committed to the truncation/rescaling asymmetry
appearing in survivorship-selected populations generally. And it would
be committed to the exponent's dependence on noise placement, $\tfrac14$
for integrated drivers, $\tfrac12$ for Brownian ones, being visible
in any domain where hazard dynamics can be measured directly. None of
the three is established here, all three could be found false without
touching a line of Section~\ref{sec:imports}, and that is the
point of naming them.


\section{The prospective test}\label{sec:protocol}

The Route World's flagship empirical discriminator is not any of the
retrospective cases above, which are labeled exploratory where they
appear, but a test declared before its observation interval. The
retrospective-target objection is answered in form by the imported
theorem, which fixes the target before the record is read and the
alternative to each attribute as its complement; what no theorem closes
is the choice of which attributes to interrogate, and only a declared
scan closes it. The test is therefore declared here, in the form the
theorem companion prescribes. The target family is survival to horizon
$t$ of the biosphere, read through the Golden target of
Section~\ref{sec:target}; the rival law is the observer-conditioned
draw. The attribute family is the set of prerequisites the Golden target
is argued to require and the observer target is insensitive to, fixed in
number before scoring, with the dependence assumption stated as
submultiplicativity or a bound on the conjunction. The baseline priors are
the Cosmic-Nexus weights the natural-history leg's method estimates. The
record statistic is the conjunction of those prerequisites observed over
the interval, and the reported quantity is the Bayes factor between the
conditioned and unconditioned laws. The stopping rule is a fixed
interval, and the multiplicity budget is the expected number of spurious
$\kappa$-fold flags at the declared $\kappa$. The provisioning test, which
asks whether departure-relevant resources sit where the Golden target
requires them at a rate above the habitability-conditioned base rate,
is the same protocol at a second grain and awaits an exoplanet census.
Until a test so declared has been scored, the Route World is a hypothesis
with an estimated Bayes factor from retrospective material, and it is
described that way.

\section{Open problems}\label{sec:open}

Every [O] tag in the paper is gathered here. The estimate of the
natural-history leg's exponent $\alpha$ by a direct attribute-level audit
of the record. The identification map from accordance-legible traces to
the target's operation. The formal dissolution of the finale paradox at
the grade of the draw. Strong Past-Alignment, conditional on the
cumulative-suppression landscape assumption, which is the Optimal-Earth
companion's and not this paper's. The landscape premise under which every
alternative lineage is an Imperfection Magnitude Drop. The shifted
convergence result on which any forward reading of practical Goldenness
depends. The provisioning base rate. An explicit intermediate law for the Weak reading, a
partial-conditioning family whose expected log-likelihood growth can be
computed and ordered between the two rigorous signs. And the Route World
itself, entered as a hypothesis: what is derived is its structure, what
is open is whether the record is under it.

\section{Falsifiers}\label{sec:falsifiers}

What would count against the hypothesis, in one place. A Bayes factor
between the conditioned and observer-conditioned laws, on a predeclared
statistic, found near or below one. The information flux measured and
found at the rate the unconditioned law predicts. The provisioning base
rate found undistinguished from the habitability-conditioned one. A
perturbation exhibited that leaves a Golden-relevant trait's Cosmic-Nexus
weight undiminished, which weakens local alignment, and a systematic class
of such perturbations, which breaks it. A route-deleting fluke whose
complement is shown at least as Golden-productive, at a rate above the
prior bound's $p$, or a stockpile whose members fail the route-deletion
audit individually or whose family fails submultiplicativity. And a
positive calibration in the natural-history leg that fails the once-only
test. Against the ceiling no falsifier is offered, because it is a
theorem; what it forbids is reading any of these, unfired, as
certification.

\section{Conclusion}\label{sec:conclusion}

The claim stated at the outset is a hypothesis and stands as one. Under
the stated counterfactual-weight and target-requirement assumptions, the
record of natural history yields a Bayes factor favoring conditioning on
the biosphere's indefinite survival over conditioning on the existence
of observers, estimated from the record's technological-route
prerequisites and bounded from its natural-historical counterfactual
weight; the scientific burden now sits entirely in establishing the two
probabilities independently of the record they are scored on, which is
what the prospective protocol is for, and until it has been run the
factor is a retrospective estimate under assumptions and is described
that way. Every consequence drawn from the hypothesis carries a test. The theorem that makes the hypothesis well posed, that a
future-targeted draw leaves fingerprints on the past without acting on
it, is the companion's, and this paper adds to it only the target, which
it constructs, the arguments, which it states as comparisons, the
instrument, which it reads as an instance, and the discriminators, which
it names. What the hypothesis cannot have is certification, by the
theorem it rests on; what it can have is a Bayes factor and a declared
test, and the posture that follows is the one the mathematics imposes:
held, not possessed.

\begin{thebibliography}{99}\small
\bibitem{Aethus} A.~Benander, \emph{Aethic Reasoning: A
Comprehensive Solution to the Quantum Measurement Problem}, November 2024;
PhilPapers BENARA-5, PhilSci-Archive eprint 25713.

\bibitem{AethusBrief} A.~Benander, \emph{Aethic Reasoning in
Brief: From a Counterfactual Semantics to the Quantum Measurement
Problem}, preprint, 2026, PhilSci-Archive eprint 30439.

\bibitem{AethusDecision} A.~Benander, \emph{Newcomb's Problem Without a
Fixed Past: A Decision Rule From Observer-Indexed Determinacy}, working
draft, 2026.

\bibitem{Allen09} R.~C.~Allen, \emph{The British Industrial
Revolution in Global Perspective}, Cambridge University Press, 2009.

\bibitem{Arrian} Arrian, \emph{Anabasis of Alexander}, 1.15.8.

\bibitem{BarrowTipler86} J.~D.~Barrow and F.~J.~Tipler, \emph{The
Anthropic Cosmological Principle}, Oxford University Press, Oxford,
1986.
\bibitem{Ben26Optimal} A.~Benander, \emph{The Optimal Earth: A
Paradigmatic Re-Rendition of Natural History from Coupled
Improbability}, working draft, 2026.

\bibitem{BenGradeSchool} A.~K.~Benander, \emph{Ajax's Math Ideas}
(grade-school idea journal), dated public record,
\url{https://aethus.org/aethic-code/history/ajax-grade-school-ideas.html};
Golden Point entry mid-January 2021.
\bibitem{ConwayMorris03} S.~Conway Morris, \emph{Life's Solution:
Inevitable Humans in a Lonely Universe}, Cambridge University Press,
2003.
\bibitem{Diamond97} J.~Diamond, \emph{Guns, Germs, and Steel: The
Fates of Human Societies}, W.~W.~Norton, 1997.

\bibitem{Eisenstein79} E.~L.~Eisenstein, \emph{The Printing Press as an
Agent of Change}, Cambridge University Press, 1979.
\bibitem{GoldenScience22} A.~Benander, \emph{Golden Science},
manuscript, August 2022; archived copy dated November 2022.
\newblock \url{https://aethus.org/aethic-code/files/GoldenScience_Nov2022.pdf}

\bibitem{Gould89} S.~J.~Gould, \emph{Wonderful Life}, W.~W.~Norton, New
York, 1989.
\bibitem{Leslie79} J.~Leslie, \emph{Value and Existence}, Basil
Blackwell, Oxford, 1979.
\bibitem{Merton38} R.~K.~Merton, \emph{Science, technology and society
in seventeenth century England}, Osiris \textbf{4} (1938), 360--632.
\bibitem{NexicFoundation} A.~Benander, \emph{The Accordance Principle:
Foundations of Nexic Reasoning}, working draft, 2026.

\bibitem{Nexus} A.~Benander, \emph{Nexic Reasoning: Defining
a Generalized Calculus Over Anthropic Parameters}, 2024,
\texttt{doi:10.31219/osf.io/jvkcp}.

\bibitem{PlutarchAlex} Plutarch, \emph{Life of Alexander},
\S16, in \emph{Parallel Lives}.

\bibitem{Pomeranz00} K.~Pomeranz, \emph{The Great Divergence:
China, Europe, and the Making of the Modern World Economy}, Princeton
University Press, 2000.

\bibitem{RussellSeguin82} D.~A.~Russell and R.~S\'eguin,
\emph{Reconstructions of the small Cretaceous theropod
Stenonychosaurus inequalis and a hypothetical dinosauroid}, Syllogeus
\textbf{37} (1982), 1--43 (AUTHOR TO VERIFY).
\bibitem{Teilhard55} P.~Teilhard de Chardin, \emph{The Phenomenon of
Man}, Harper, New York, 1959 (French original 1955).
\bibitem{Tipler94} F.~J.~Tipler, \emph{The Physics of Immortality},
Doubleday, New York, 1994.
\bibitem{WeightAlg} A.~Benander, \emph{The Weight Algebra of
Information States: A Two-Layer Commutative Semiring for Superposition
and Pruning}, working draft, July 2026.

\bibitem{SurvA} A.~Benander, \emph{Survivorship conditioning for hazard rates driven at the log-derivative level}, working draft, 2026.
\bibitem{TargetedB1} A.~Benander, \emph{Conditioning a record on its own future: targeted accordance and survivorship-conditioned path measures}, working draft, 2026.
\bibitem{TechnoC} A.~Benander, \emph{Technoanatomy: civilizational survival as a sign, not a size}, working draft, 2026.

\end{thebibliography}
\end{document}
